% આ ભાગનો પૂર્ણ સંપાદનીય પાઠ છે. શૈલી, ફોન્ટ, ચિત્રો અને પ્રકરણસંદર્ભ માટે સંપૂર્ણ સ્રોત ZIP તથા reader/full-reader.tex જુઓ. આ એકલી સ્વતંત્ર કમ્પાઇલ થતી મુખ્ય ફાઇલ નથી.

% BEGIN OLP-0322
\olpart{sol}{દ્વિતીય-ક્રમ તર્કશાસ્ત્ર}
\phantomsection\label{guunit:OLP-0322}


\begin{editorial}
આ દ્વિતીય-ક્રમ તર્કશાસ્ત્રના ભાગની શરૂઆત છે.
\end{editorial}

\begingroup
\makeatletter\def\ol@sectioncs{section}\makeatother

% BEGIN OLP-0323
\olchapter{sol}{syn}{વાક્યરચના અને અર્થવિજ્ઞાન}
\olchapterid{sol}{syn}
\phantomsection\label{guunit:OLP-0323}


\begin{editorial}
હાલમાં આ પ્રકરણમાં દ્વિતીય-ક્રમ તર્કશાસ્ત્રની મૂળભૂત વાક્યરચના અને
અર્થવિજ્ઞાન આવરી લેવાયાં છે. પ્રકરણ તરીકે તે હજી ઘણું ટૂંકું છે.
દ્વિતીય-ક્રમ તર્કશાસ્ત્રની નિષ્પત્તિ પદ્ધતિઓની ચર્ચા કરી શકાય તે
માટે દ્વિતીય-ક્રમ ચલોનું પ્રતિસ્થાપન પણ સમજાવવું પડશે; તેમાં કેટલીક
સૂક્ષ્મ મુશ્કેલીઓ છે.
\end{editorial}

\begingroup
\makeatletter\def\ol@sectioncs{section}\makeatother

% BEGIN OLP-0324
\olfileid{sol}{syn}{int}

\olsection{પ્રસ્તાવના}
\phantomsection\label{guunit:OLP-0324}


પ્રથમ-ક્રમ તર્કશાસ્ત્રમાં આપેલી ભાષાનાં અતાર્કિક સંકેતો, એટલે કે તેનાં
અચળો, ફલનો અને વિધેયો, તાર્કિક સંકેતો સાથે
જોડીને આપણે પ્રથમ-ક્રમ સંરચનાઓ વિશેનાં વિધાનો વ્યક્ત કરીએ છીએ.
આ માટે સંતોષનો ખ્યાલ વપરાય છે. તે સંરચના~$\Struct{M}$,
ચલ-નિયુક્તિ~$s$ અને સૂત્ર~$!A$ને જોડે છે: $\Sat{M}{!A}[s]$
ત્યારે અને માત્ર ત્યારે જ સાચું છે જ્યારે~$!A$નાં અચળો,
ફલનો અને વિધેયોનું~$\Struct{M}$ મુજબ અર્થઘટન તથા
તેના મુક્ત ચલોનું~$s$ મુજબ અર્થઘટન કરીએ ત્યારે તે જે કહે છે તે
સાચું હોય. તાદાત્મ્ય~$\eq$નું અર્થઘટન, તેમજ~$\lforall$ અને
$\lexists$નું અર્થઘટન, $\Sat{M}{!A}[s]$ની વ્યાખ્યામાં જ નિહિત છે.
સમાનતા હંમેશાં સંરચનાના પ્રદેશ~$\Domain{M}$ પરના અભિન્નતા
સંબંધ તરીકે અર્થઘટિત થાય છે અને પરિમાણકો હંમેશાં આખા પ્રદેશ પર
વ્યાપે છે. પરંતુ અગત્યની મર્યાદા એ છે કે પરિમાણીકરણ માત્ર
પ્રદેશના ઘટકો પર જ થઈ શકે છે; તેથી પરિમાણક પછી માત્ર વસ્તુ-ચલો
આવી શકે છે.

દ્વિતીય-ક્રમ તર્કશાસ્ત્રમાં ભાષા અને સંતોષની વ્યાખ્યા બંનેને મુક્ત અને
બદ્ધ વિધેય-ચલો તથા સંબંધ-ચલો અને તેમના પરના પરિમાણીકરણનો સમાવેશ થાય
એ રીતે વિસ્તૃત કરીએ છીએ. વસ્તુ-ચલોનો અચળો સાથે જેવો સંબંધ છે,
એવો જ આ ચલોનો ફલનો અને વિધેયો સાથે છે. દ્વિતીય-ક્રમ
પદો અને સૂત્રોની રચનામાં તેઓ એ જ પ્રકારની ભૂમિકા ભજવે છે અને
તેમના પરનું પરિમાણીકરણ પણ સમાન રીતે સંભાળવામાં આવે છે. \emph{પ્રમાણભૂત}
અર્થવિજ્ઞાનમાં દ્વિતીય-ક્રમ પરિમાણકો યોગ્ય પ્રકારની બધી શક્ય વસ્તુઓ
પર વ્યાપે છે: વિધેય-ચલો માટે $\Domain{M}$માંથી~$\Domain{M}$માં જતાં
$n$-સ્થાની વિધેયો અને વિધેયધર્મ-ચલો માટે $n$-સ્થાની સંબંધો.
દાખલા તરીકે,
$\lforall[\Obj{v_0}][(\Obj{P^1_0}(\Obj{v_0}) \lor \lnot
  \Obj{P^1_0}(\Obj{v_0}))]$ પ્રથમ-ક્રમ અને દ્વિતીય-ક્રમ બંને
તર્કશાસ્ત્રમાં સૂત્ર છે. દ્વિતીય-ક્રમ તર્કશાસ્ત્રમાં વળી
$\lforall[\Obj{V^1_0}][\lforall[\Obj{v_0}][(\Obj{V^1_0}(\Obj{v_0}) \lor
    \lnot \Obj{V^1_0}(\Obj{v_0}))]]$ અને
$\lexists[\Obj{V^1_0}][\lforall[\Obj{v_0}][(\Obj{V^1_0}(\Obj{v_0}) \lor
    \lnot \Obj{V^1_0}(\Obj{v_0}))]]$ પણ વિચારી શકાય. તેમાં કોઈ મુક્ત
ચલ નથી, તેથી તે દ્વિતીય-ક્રમ તર્કશાસ્ત્રનાં વાક્યો છે. અહીં
$\Obj{V^1_0}$ દ્વિતીય-ક્રમનો $1$-સ્થાની વિધેયધર્મ-ચલ છે. $\Obj{V^1_0}$નું
સંભવિત અર્થઘટન~$\Obj{P^1_0}$ જેવા $1$-સ્થાની વિધેયને આપી
શકાતા અર્થઘટન જેવું જ છે, એટલે~$\Domain{M}$ના ઉપગણો. તેથી આવા ચલો
પર પરિમાણીકરણનો અર્થ એ થાય છે કે
$\lforall[\Obj{v_0}][(\Obj{V^1_0}(\Obj{v_0}) \lor \lnot
  \Obj{V^1_0}(\Obj{v_0}))]$ એ~$\Obj{V^1_0}$ને~$\Domain{M}$નો કોઈપણ
ઉપગણ આપીએ ત્યારે સાચું હોય છે, અથવા ઓછામાં ઓછો એક એવો ઉપગણ આપીએ
ત્યારે સાચું હોય છે. દરેક ગણ આપેલા ઘટકને કાં તો સમાવે છે અથવા નથી
સમાવતો, તેથી બંને દ્વિતીય-ક્રમ વાક્યો કોઈપણ સંરચનામાં સાચાં છે.
\sourcecorrection{OLSOL-001}{સ્થિર અંગ્રેજી મૂળના છેલ્લા સૂત્રમાં
$\Obj v_0$ અને પછી $v_0$ લખાયું છે. ઉપરના બદ્ધ વસ્તુ-ચલ માટેના
એકસરખા $\Obj{v_0}$ સંકેતથી બંને સ્થાનો સ્પષ્ટ કર્યાં છે.}
% END OLP-0324
\endgroup


\begingroup
\makeatletter\def\ol@sectioncs{section}\makeatother

% BEGIN OLP-0325
\olfileid{sol}{syn}{frm}

\olsection{પદો અને સૂત્રો}
\phantomsection\label{guunit:OLP-0325}


પ્રથમ-ક્રમ તર્કશાસ્ત્રની જેમ જ દ્વિતીય-ક્રમ તર્કશાસ્ત્રની અભિવ્યક્તિઓ
પાયાના એવા સંકેતભંડોળમાંથી બને છે જેમાં \emph{ચલો},
\emph{અચળો}, \emph{વિધેયો} અને ક્યારેક
\emph{ફલનો} હોય છે. તેમની સાથે તાર્કિક સંયોજકો, પરિમાણકો
અને કૌંસ તથા અલ્પવિરામ જેવા વિરામસંકેતો જોડીને \emph{પદો} અને
\emph{સૂત્રો} રચાય છે. ફરક એ છે કે વસ્તુઓ માટેના ચલો ઉપરાંત
દ્વિતીય-ક્રમ તર્કશાસ્ત્રમાં સંબંધો અને વિધેયો માટેના ચલો પણ હોય છે,
અને તેમના પર પરિમાણીકરણ કરી શકાય છે. આમ, દ્વિતીય-ક્રમ તર્કશાસ્ત્રના
તાર્કિક સંકેતોમાં પ્રથમ-ક્રમ તર્કશાસ્ત્રના સંકેતો ઉપરાંત આનો સમાવેશ થાય છે:

\begin{enumerate}
\item દરેક સ્થાનસંખ્યા~$n$ માટે દ્વિતીય-ક્રમ સંબંધ-ચલોનો
  ગણનીય ગણ: $\Obj V_0^n$, $\Obj V_1^n$, $\Obj V_2^n$, \dots
\item દ્વિતીય-ક્રમ વિધેય-ચલોનો ગણનીય ગણ:
  $\Obj u_0^n$, $\Obj u_1^n$, $\Obj u_2^n$, \dots
\end{enumerate}

જેમ આપણે પ્રથમ-ક્રમ ચલો~$\Obj v_i$ માટે $x$, $y$, $z$ને પરાચલો
તરીકે વાપરીએ છીએ, તેમ~$\Obj V_i^n$ માટે $X$, $Y$, $Z$ વગેરે અને
$\Obj u_i^n$ માટે $u$, $v$ વગેરે પરાચલો તરીકે વાપરીશું.

\begin{explain}
દ્વિતીય-ક્રમ ભાષાના અતાર્કિક સંકેતો પ્રથમ-ક્રમ ભાષાની જેમ જ, તેના
અચળો, ફલનો અને વિધેયોની યાદી દ્વારા
નિર્દિષ્ટ થાય છે.

પ્રથમ-ક્રમ તર્કશાસ્ત્રમાં સામાન્ય રીતે તાદાત્મ્ય~$\eq$નો પણ સમાવેશ
કરવામાં આવે છે. પ્રથમ-ક્રમ તર્કશાસ્ત્રમાં કંઈ રસપ્રદ વ્યક્ત કરી શકાય
તે માટે ભાષા~$\Lang{L}$ના અતાર્કિક સંકેતો બહુ મહત્વના છે. અલબત્ત,
અતાર્કિક સંકેત વિના પણ કેટલાંક વાક્યો હોય છે, પરંતુ માત્ર~$\eq$
વડે કંઈ રસપ્રદ કહેવું મુશ્કેલ છે. દ્વિતીય-ક્રમ તર્કશાસ્ત્રમાં સંબંધો
અને વિધેયો માટે અનિયંત્રિત સંખ્યામાં ચલો હોવાથી, અતાર્કિક સંકેતોનો
વિશેષ ભંડાર ન હોય તો પણ પ્રથમ-ક્રમ ભાષામાં કહી શકાય એવી વાતો કહી
શકાય છે.
\end{explain}

\begin{defn}[દ્વિતીય-ક્રમ પદો]
ભાષા~$\Lang L$નાં \emph{દ્વિતીય-ક્રમ પદો}નો ગણ~$\TrmSOL[L]$
\olref[fol][syn][frm]{defn:terms}માં નીચેની કલમ ઉમેરવાથી વ્યાખ્યાયિત
થાય છે:
\begin{enumerate}
\item જો $u$ કોઈ $n$-સ્થાની વિધેય-ચલ હોય અને $t_1$, \dots, $t_n$
  પદો હોય, તો $\Atom{u}{t_1, \ldots, t_n}$ પદ છે.
\end{enumerate}
\end{defn}

\begin{explain}
આમ, દ્વિતીય-ક્રમ પદ પ્રથમ-ક્રમ પદ જેવું જ હોય છે, સિવાય કે પ્રથમ-ક્રમ
પદમાં જ્યાં ફલન~$\Obj{f^n_i}$ હોય ત્યાં દ્વિતીય-ક્રમ પદમાં
તેના સ્થાને વિધેય-ચલ~$\Obj{u^n_i}$ હોઈ શકે.
\end{explain}

\begin{defn}[દ્વિતીય-ક્રમ સૂત્ર]
ભાષા~$\Lang L$નાં \emph{દ્વિતીય-ક્રમ સૂત્રો}નો
ગણ~$\FrmSOL[L]$, \olref[fol][syn][frm]{defn:formulas}માં નીચેની કલમો
ઉમેરવાથી વ્યાખ્યાયિત થાય છે:
\begin{enumerate}
\item જો $X$ કોઈ $n$-સ્થાની વિધેયધર્મ-ચલ હોય અને $t_1$, \dots,
  $t_n$ ભાષા~$\Lang L$નાં દ્વિતીય-ક્રમ પદો હોય, તો
  $\Atom{X}{t_1,\ldots, t_n}$ આણ્વિક સૂત્ર છે.

\item જો $!A$ સૂત્ર હોય અને $u$ વિધેય-ચલ હોય,
  તો $\lforall[u][!A]$ સૂત્ર છે.

\item જો $!A$ સૂત્ર હોય અને $X$ વિધેયધર્મ-ચલ હોય,
  તો $\lforall[X][!A]$ સૂત્ર છે.

\item જો $!A$ સૂત્ર હોય અને $u$ વિધેય-ચલ હોય,
  તો $\lexists[u][!A]$ સૂત્ર છે.

\item જો $!A$ સૂત્ર હોય અને $X$ વિધેયધર્મ-ચલ હોય,
  તો $\lexists[X][!A]$ સૂત્ર છે.
\end{enumerate}
\end{defn}
% END OLP-0325
\endgroup


\begingroup
\makeatletter\def\ol@sectioncs{section}\makeatother

% BEGIN OLP-0326
\olfileid{sol}{syn}{sat}

\olsection{સંતોષ}
\phantomsection\label{guunit:OLP-0326}


\begin{explain}
દ્વિતીય-ક્રમ સૂત્રો માટે સંતોષ-સંબંધ~$\Sat{M}{!A}[s]$
વ્યાખ્યાયિત કરવા માટે આપણે દ્વિતીય-ક્રમ ચલોને પણ આવરી લે
એ રીતે અગાઉની વ્યાખ્યાઓ વિસ્તૃત કરવી પડે. સંરચનાનો ખ્યાલ
દ્વિતીય-ક્રમ અને પ્રથમ-ક્રમ તર્કશાસ્ત્રમાં એકસરખો છે. ફરક માત્ર
ચલ-નિયુક્તિઓ~$s$માં છે: હવે તેમણે પ્રથમ-ક્રમ ચલો ઉપરાંત
દ્વિતીય-ક્રમ ચલોને પણ મૂલ્યો આપવા પડે.
\end{explain}

\begin{defn}[ચલ-નિયુક્તિ]
સંરચના~$\Struct{M}$ માટેની \emph{ચલ-નિયુક્તિ}~$s$ એવું
વિધેય છે જે દરેક
\begin{enumerate}
\item વસ્તુ-ચલ~$\Obj{v_i}$ને~$\Domain M$નો ઘટક આપે છે,
  એટલે $s(\Obj{v_i}) \in \Domain{M}$;
\item $n$-સ્થાની સંબંધ-ચલ~$\Obj{V_i^n}$ને~$\Domain{M}$ પરનો
  $n$-સ્થાની સંબંધ આપે છે, એટલે
  $s(\Obj{V_i^n}) \subseteq \Domain{M}^n$;
\item $n$-સ્થાની વિધેય-ચલ~$\Obj{u_i^n}$ને $\Domain{M}$માંથી
  $\Domain{M}$માં જતું $n$-સ્થાની વિધેય આપે છે, એટલે
  $s(\Obj{u_i^n})\colon \Domain{M}^n \to \Domain{M}$.
\end{enumerate}
\end{defn}

\begin{explain}
સંરચના દરેક અચળ અને ફલનને મૂલ્ય
આપે છે; દ્વિતીય-ક્રમ ચલ-નિયુક્તિ દરેક વસ્તુ-ચલને વસ્તુ, દરેક
વિધેય-ચલને વિધેય અને દરેક સંબંધ-ચલને સંબંધ આપે છે. આ બંને સાથે
મળીને દરેક પદનું મૂલ્ય નક્કી કરી શકાય છે.
\sourcecorrection{OLSOL-002}{સ્થિર અંગ્રેજી મૂળની આ સમજણમાં
ચલ-નિયુક્તિ દ્વારા વસ્તુ-ચલો અને વિધેય-ચલોને મળતા મૂલ્યો જણાવ્યાં
છે, પરંતુ ઉપરની વ્યાખ્યામાં આવરી લીધેલા સંબંધ-ચલોને મળતા સંબંધોનો
ઉલ્લેખ રહી ગયો છે. અહીં તે ત્રીજો પ્રકાર પણ સ્પષ્ટ કર્યો છે.}
\end{explain}

\begin{defn}[પદનું મૂલ્ય]
જો~$t$ ભાષા~$\Lang L$નું પદ હોય, $\Struct M$ એ~$\Lang L$ માટેની
સંરચના હોય અને~$s$ એ~$\Struct M$ માટેની ચલ
નિયુક્તિ હોય, તો \emph{મૂલ્ય}~$\Value{t}{M}[s]$ પ્રથમ-ક્રમ
પદોની જેમ જ, નીચેની વધારાની કલમ સાથે વ્યાખ્યાયિત થાય છે:
\begin{quote}
\indcase{t}{\Atom{u}{t_1, \ldots, t_n}}{
\[
\Value{\indfrm}{M}[s] = s(u)(\Value{t_1}{M}[s], \ldots,
\Value{t_n}{M}[s]).
\]}
\end{quote}
\end{defn}

\begin{defn}[$x$-રૂપાંતર]
જો~$s$ એ સંરચના~$\Struct M$ માટેની ચલ નિયુક્તિ
હોય, તો~$\Struct M$ માટેની એવી કોઈપણ ચલ નિયુક્તિ~$s'$,
જે~$s$થી વધુમાં વધુ~$x$ને આપેલા મૂલ્યમાં જ જુદી હોય, તેને~$s$નું
\emph{$x$-રૂપાંતર} કહે છે. $s'$ એ~$s$નું $x$-રૂપાંતર હોય તો
$\varAssign{s'}{s}{x}$ લખીએ છીએ. (દ્વિતીય-ક્રમ ચલો~$X$ અથવા~$u$
માટે પણ આવું જ છે.)
\end{defn}

\begin{defn}
  જો~$s$ એ સંરચના~$\Struct M$ માટેની ચલ નિયુક્તિ
  હોય અને~$m \in \Domain{M}$ હોય, તો~$\Subst{s}{m}{x}$ નીચે મુજબ
  વ્યાખ્યાયિત ચલ નિયુક્તિ છે:
  \[\Subst{s}{m}{y} = \begin{cases}
    m & \text{જો } y \ident x\\
    s(y) & \text{અન્યથા},
  \end{cases}\]
  જો~$X$ એ $n$-સ્થાની સંબંધ-ચલ હોય અને $M \subseteq
  \Domain{M}^n$ હોય, તો~$\Subst{s}{M}{X}$ નીચે મુજબ વ્યાખ્યાયિત
  ચલ નિયુક્તિ છે:
  \[\Subst{s}{M}{y} = \begin{cases}
    M & \text{જો } y \ident X\\
    s(y) & \text{અન્યથા}.
  \end{cases}\]
  જો~$u$ એ $n$-સ્થાની વિધેય-ચલ હોય અને $f\colon \Domain{M}^n
  \to \Domain{M}$ હોય, તો~$\Subst{s}{f}{u}$ નીચે મુજબ વ્યાખ્યાયિત
  ચલ નિયુક્તિ છે:
  \[\Subst{s}{f}{y} = \begin{cases}
    f & \text{જો } y \ident u\\
    s(y) & \text{અન્યથા}.
  \end{cases}\]
  દરેક કિસ્સામાં~$y$ કોઈપણ પ્રથમ-ક્રમ અથવા દ્વિતીય-ક્રમ ચલ
  હોઈ શકે છે.
\end{defn}

\begin{defn}[સંતોષ]
દ્વિતીય-ક્રમ સૂત્રો~$!A$ માટે સંતોષની વ્યાખ્યા
\olref[fol][syn][sat]{defn:satisfaction} જેવી જ છે, પરંતુ તેમાં
નીચેની કલમો ઉમેરાય છે:
\begin{enumerate}
\item \indcase{!A}{\Atom{X^n}{t_1, \dots, t_n}}{$\Sat{M}{\indfrm}[s]$
  ત્યારે અને માત્ર ત્યારે જ જ્યારે
  $\langle \Value{t_1}{M}[s], \dots, \Value{t_n}{M}[s] \rangle \in
  s(X^n)$.}
\item %
  \indcase{!A}{\lforall[X][!B]}{$\Sat{M}{\indfrm}[s]$ ત્યારે અને
  માત્ર ત્યારે જ જ્યારે દરેક $M \subseteq \Domain{M}^n$ માટે
  $\Sat{M}{!B}[\Subst{s}{M}{X}]$.}

\item %
  \indcase{!A}{\lexists[X][!B]}{$\Sat{M}{\indfrm}[s]$ ત્યારે અને
  માત્ર ત્યારે જ જ્યારે ઓછામાં ઓછા એક $M \subseteq \Domain{M}^n$
  માટે $\Sat{M}{!B}[\Subst{s}{M}{X}]$.}

\item %
  \indcase{!A}{\lforall[u][!B]}{$\Sat{M}{\indfrm}[s]$ ત્યારે અને
  માત્ર ત્યારે જ જ્યારે દરેક
  $f\colon \Domain{M}^n \to \Domain{M}$ માટે
  $\Sat{M}{!B}[\Subst{s}{f}{u}]$.}

\item %
  \indcase{!A}{\lexists[u][!B]}{$\Sat{M}{\indfrm}[s]$ ત્યારે અને
  માત્ર ત્યારે જ જ્યારે ઓછામાં ઓછા એક
  $f\colon \Domain{M}^n \to \Domain{M}$ માટે
  $\Sat{M}{!B}[\Subst{s}{f}{u}]$.}
\end{enumerate}
\end{defn}

\begin{ex}
  સૂત્ર~$\lforall[z][(\Atom{X}{z} \liff \lnot
    \Atom{Y}{z})]$ વિચારો. તેમાં દ્વિતીય-ક્રમ પરિમાણકો નથી, પરંતુ
  દ્વિતીય-ક્રમ ચલો~$X$ અને~$Y$ છે (અહીં બંને એકસ્થાની
  સમજવાના છે). તેને અનુરૂપ પ્રથમ-ક્રમ વાક્ય
  $\lforall[z][(\Atom{P}{z} \liff \lnot \Atom{R}{z})]$ કહે છે કે
  જે કશું~$P$ના અર્થઘટનમાં આવે છે તે~$R$ના અર્થઘટનમાં આવતું નથી,
  અને ઊલટું પણ. કોઈ સંરચનામાં વિધેય~$P$નું અર્થઘટન
  $\Assign{P}{M}$ વડે મળે છે. પરંતુ~$X$ અને~$Y$ જેવા દ્વિતીય-ક્રમ
  ચલોનું અર્થઘટન સંરચના પોતે નહિ, પરંતુ
  ચલ નિયુક્તિ આપે છે. આ દ્વિતીય-ક્રમ સૂત્ર,
  વાક્ય નથી (તેમાં મુક્ત ચલો~$X$ અને~$Y$ છે), તેથી
  તેનો સંતોષ સંરચના~$\Struct{M}$ અને તેની સાથેની
  ચલ નિયુક્તિ~$s$ને સાપેક્ષ જ નક્કી થાય છે.

  વ્યાપકક્ષેત્રના બધા ઘટકોને ધ્યાનમાં લઈએ તો, $s(X)$ના ઘટકો
  ચોક્કસ એ જ છે જે $s(Y)$ના ઘટકો નથી; એટલે દરેક ઘટક માટે આ
  દ્વિમુખી શરત સાચી હોય,
  ત્યારે $\Sat{M}{\lforall[z][(\Atom{X}{z} \liff \lnot
  \Atom{Y}{z})]}[s]$ સાચું છે; એટલે તે ત્યારે અને માત્ર ત્યારે જ
  જ્યારે $s(Y) = \Domain{M} \setminus s(X)$.
  \sourcecorrection{OLSOL-003}{સ્થિર અંગ્રેજી મૂળ અહીં બંને ગણોના
  ઘટકો એકબીજામાં નથી એવું કહે છે, જે માત્ર વિચ્છેદતા બતાવે છે;
  સૂત્રની દ્વિમુખી શરત મુજબ બંને ગણ પરસ્પર પૂરક હોવા જરૂરી છે.
  અહીં દરેક ઘટક માટે તે શરત સ્પષ્ટ કરી છે.}
  દાખલા તરીકે,
  $\Domain{M} = \{1, 2, 3\}$ લો. અહીં કોઈ વિધેયો,
  ફલનો કે અચળો આવતાં નથી, એટલે
  $\Struct{M}$નું વ્યાપકક્ષેત્ર જ સંબંધિત છે. હવે
  $s_1(X) = \{1, 2\}$ અને $s_1(Y) = \{3\}$ માટે
  $\Sat{M}{\lforall[z][(\Atom{X}{z}
      \liff \lnot \Atom{Y}{z})]}[s_1]$ મળે છે.

  બીજી તરફ, જો $s_2(X) = \{1, 2\}$ અને $s_2(Y) = \{2, 3\}$ હોય,
  તો $\Sat/{M}{\lforall[z][(\Atom{X}{z} \liff \lnot
      \Atom{Y}{z})]}[s_2]$. કારણ કે
  $\Sat{M}{\Atom{X}{z}}[\Subst{s_2}{2}{z}]$ ($2 \in \Subst{s_2}{2}{z}(X)$ હોવાથી),
  પરંતુ $\Sat/{M}{\lnot \Atom{Y}{z}}[\Subst{s_2}{2}{z}]$
  (કારણ કે $2 \in \Subst{s_2}{2}{z}(Y)$ પણ છે).
\end{ex}

\begin{ex}
  કોઈ $N \subseteq \Domain{M}$ માટે
  $\Sat{M}{(\lexists[y][\Atom{Y}{y}] \land \lforall[z][(\Atom{X}{z}
  \liff \lnot \Atom{Y}{z})])}[\Subst{s}{N}{Y}]$ હોય, તો
  $\Sat{M}{\lexists[Y][(\lexists[y][\Atom{Y}{y}] \land
  \lforall[z][(\Atom{X}{z} \liff \lnot \Atom{Y}{z})])]}[s]$ હોય છે.
  આ માટે $N \neq \emptyset$ હોવું પડે (જેથી
  $\Sat{M}{\lexists[y][\Atom{Y}{y}]}[\Subst{s}{N}{Y}]$ થાય) અને
  અગાઉના ઉદાહરણની જેમ $N = \Domain{M} \setminus s(X)$ હોવું પડે.
  બીજા શબ્દોમાં,
  $\Sat{M}{\lexists[Y][(\lexists[y][\Atom{Y}{y}] \land
  \lforall[z][(\Atom{X}{z} \liff \lnot \Atom{Y}{z})])]}[s]$
  ત્યારે અને માત્ર ત્યારે જ જ્યારે $\Domain{M} \setminus s(X)$
  અરિક્ત હોય, એટલે $s(X) \neq \Domain{M}$. આમ,
  $\Domain{M} = \{1, 2, 3\}$ અને $s(X) = \{1, 2\}$ હોય તો આ
  સૂત્ર સંતોષાય છે, પરંતુ $s(X) = \{1, 2, 3\} = \Domain{M}$
  હોય તો સંતોષાતું નથી.

  $s(X) = \Domain{M}$ હોય ત્યારે આ સૂત્ર સંતોષાતું નથી, તેથી
  વાક્ય
  \[
  \lforall[X][\lexists[Y][(\lexists[y][\Atom{Y}{y}] \land
      \lforall[z][(\Atom{X}{z} \liff \lnot \Atom{Y}{z})])]]
  \]
  ક્યારેય સંતોષાતું નથી: કોઈપણ સંરચના~$\Struct{M}$ માટે
  $s(X) = \Domain{M}$ એવી નિયુક્તિ વાક્યને અસત્ય બનાવશે.
  બીજી બાજુ, વાક્ય
  \[
  \lexists[X][\lexists[Y][(\lexists[y][\Atom{Y}{y}] \land
      \lforall[z][(\Atom{X}{z} \liff \lnot \Atom{Y}{z})])]]
  \]
  કોઈપણ નિયુક્તિ~$s$ને સાપેક્ષ સંતોષાય છે, કારણ કે હંમેશાં
  $M \subseteq \Domain{M}$ અને $M \neq \Domain{M}$ એવો ઉપગણ મળી
  શકે છે (દાખલા તરીકે $M = \emptyset$).
\end{ex}

\begin{ex}
  દ્વિતીય-ક્રમ વાક્ય~$\lforall[X][\lforall[y][X(y)]]$ કહે છે
  કે દરેક $1$-સ્થાની સંબંધ, એટલે દરેક ગુણધર્મ, દરેક વસ્તુને લાગુ પડે
  છે. તે સ્પષ્ટપણે ક્યારેય સાચું નથી: દરેક~$\Struct{M}$માં
  $s(X) = \emptyset$ અને $s(y) = a \in \Domain{M}$ એવી
  ચલ-નિયુક્તિ~$s$ માટે $\Sat/{M}{X(y)}[s]$ છે. આથી દ્વિતીય-ક્રમ
  તર્કશાસ્ત્રમાં $!A \lif \lforall[X][\lforall[y][X(y)]]$ એ
  $\lnot !A$ને સમકક્ષ છે, એટલે
  $\Sat{M}{!A \lif
  \lforall[X][\lforall[y][X(y)]]}$ ત્યારે અને માત્ર ત્યારે જ
  જ્યારે $\Sat{M}{\lnot !A}$. બીજા શબ્દોમાં, દ્વિતીય-ક્રમ
  તર્કશાસ્ત્રમાં $\lforall$ અને~$\lif$ વડે $\lnot$ વ્યાખ્યાયિત
  કરી શકાય છે.
\end{ex}

\begin{prob}
  બતાવો કે દ્વિતીય-ક્રમ તર્કશાસ્ત્રમાં $\lforall$ અને~$\lif$ વડે
  બીજા સંયોજકો વ્યાખ્યાયિત કરી શકાય છે:
  \begin{enumerate}
    \item સાબિત કરો કે દ્વિતીય-ક્રમ તર્કશાસ્ત્રમાં $!A \land !B$ એ
    $\lforall[X][(!A \lif (!B \lif \lforall[x][X(x)])\lif
    \lforall[x][X(x)])]$ને સમકક્ષ છે.
    \item માત્ર $\lforall$ અને $\lif$ વાપરીને $!A \lor !B$ને
    સમકક્ષ દ્વિતીય-ક્રમ સૂત્ર શોધો.
  \end{enumerate}
\end{prob}
% END OLP-0326
\endgroup


\begingroup
\makeatletter\def\ol@sectioncs{section}\makeatother

% BEGIN OLP-0327
\olfileid{sol}{syn}{sem}
\olsection{અર્થવિચારી ખ્યાલો}
\phantomsection\label{guunit:OLP-0327}


\begin{explain}
\emph{પ્રામાણ્ય}, \emph{ફલિતતા} અને \emph{સંતોષ્યતા}ના કેન્દ્રીય
તાર્કિક ખ્યાલો દ્વિતીય-ક્રમ તર્કશાસ્ત્રમાં પ્રથમ-ક્રમ તર્કશાસ્ત્રની
જેમ જ વ્યાખ્યાયિત થાય છે; માત્ર તેમનો પાયાનો સંતોષ-સંબંધ હવે
દ્વિતીય-ક્રમ સૂત્રો માટેનો છે. અલબત્ત, દ્વિતીય-ક્રમ
વાક્ય એવું સૂત્ર છે જેમાં વિધેયધર્મ-ચલો અને
વિધેય-ચલો સહિત બધા ચલો બદ્ધ હોય છે.
\end{explain}

\begin{defn}[પ્રામાણ્ય]
વાક્ય~$!A$ \emph{પ્રમાણભૂત}, $\Entails !A$, ત્યારે અને માત્ર ત્યારે જ
જ્યારે દરેક સંરચના~$\Struct M$ માટે $\Sat{M}{!A}$.
\end{defn}

\begin{defn}[ફલિતતા]
વાક્યોનો ગણ~$\Gamma$, વાક્ય~$!A$ને \emph{ફલિત કરે} છે,
$\Gamma \Entails !A$, ત્યારે અને માત્ર ત્યારે જ જ્યારે
$\Sat{M}{\Gamma}$ હોય એવી દરેક સંરચના~$\Struct M$ માટે
$\Sat{M}{!A}$.
\end{defn}

\begin{defn}[સંતોષ્યતા]
વાક્યોનો ગણ~$\Gamma$ \emph{સંતોષ્ય} છે જો કોઈ સંરચના~$\Struct M$
માટે $\Sat{M}{\Gamma}$. જો~$\Gamma$ સંતોષ્ય ન હોય, તો તેને
\emph{અસંતોષ્ય} કહે છે.
\end{defn}
% END OLP-0327
\endgroup


\begingroup
\makeatletter\def\ol@sectioncs{section}\makeatother

% BEGIN OLP-0328
\olfileid{sol}{syn}{exp}
\olsection{અભિવ્યક્તિની શક્તિ}
\phantomsection\label{guunit:OLP-0328}


\begin{explain}
દ્વિતીય-ક્રમ ચલો પરનું પરિમાણીકરણ ભાષાની અભિવ્યક્તિની શક્તિને
પ્રથમ-ક્રમ તર્કશાસ્ત્રની સરખામણીએ ઘણું વધારે છે. દ્વિતીય-ક્રમ
અસ્તિત્વવાચી પરિમાણીકરણથી અમુક ગુણધર્મો ધરાવતાં વિધેયો કે સંબંધો
અસ્તિત્વમાં છે એવું કહી શકાય છે. પ્રથમ-ક્રમ તર્કશાસ્ત્રમાં એવું કહેવા
માટે તે હેતુનો અતાર્કિક સંકેત, એટલે ફલન અથવા વિધેય,
નિર્દિષ્ટ કરવો પડે છે. દ્વિતીય-ક્રમ સર્વવાચી પરિમાણીકરણથી વ્યાપકક્ષેત્રના
બધા ઉપગણો, તેના પરના બધા સંબંધો અથવા પ્રદેશમાંથી પ્રદેશમાં
જતાં બધાં વિધેયો અમુક ગુણધર્મ ધરાવે છે એવું કહી શકાય છે. પ્રથમ-ક્રમ
તર્કશાસ્ત્રમાં માત્ર ભાષાના કોઈ અતાર્કિક સંકેતને અપાયેલા ઉપગણો,
સંબંધો કે વિધેયો તે ગુણધર્મ ધરાવે છે એવું કહી શકાય. વળી,
અમુક ગુણધર્મવાળા ઉપગણો, સંબંધો કે વિધેયો અસ્તિત્વમાં છે અથવા તે બધા
એ ગુણધર્મ ધરાવે છે એમ કહીએ ત્યારે તે ગુણધર્મ નિર્દિષ્ટ કરવામાં પણ
દ્વિતીય-ક્રમ પરિમાણીકરણ વાપરી શકાય છે. તેથી પ્રથમ-ક્રમ તર્કશાસ્ત્રમાં
અવ્યાખ્યેય સંબંધો વ્યાખ્યાયિત કરી શકાય છે અને વ્યાપકક્ષેત્રના એવા
ગુણધર્મો વ્યક્ત કરી શકાય છે જે પ્રથમ-ક્રમ તર્કશાસ્ત્રમાં વ્યક્ત થતા
નથી.
\end{explain}

\begin{defn}
જો~$\Struct{M}$ એ ભાષા~$\Lang{L}$ માટેની સંરચના હોય, તો
સંબંધ~$R \subseteq \Domain{M}^2$ એ~$\Lang{L}$માં \emph{વ્યાખ્યેય}
છે જો એવું કોઈ સૂત્ર~$!A_R(x, y)$ હોય જેમાં માત્ર $x$ અને~$y$
મુક્ત હોય અને $s(x) = a$ તથા $s(y) = b$ માટે $R(a, b)$, એટલે
$\tuple{a,b} \in R$, ત્યારે અને માત્ર ત્યારે જ જ્યારે
$\Sat{M}{!A_R(x, y)}[s]$.
\end{defn}

\begin{ex}
પ્રથમ-ક્રમ તર્કશાસ્ત્રમાં અભિન્નતા સંબંધ~$\Id{\Domain{M}}$, એટલે
$\Setabs{\tuple{a,a}}{a \in
  \Domain{M}}$, સૂત્ર~$\eq[x][y]$ વડે વ્યાખ્યાયિત કરી શકાય છે.
દ્વિતીય-ક્રમ તર્કશાસ્ત્રમાં આ સંબંધ~\emph{$\eq$ વિના} વ્યાખ્યાયિત
કરી શકાય છે. જો~$a$ અને~$b$ એ~$\Domain{M}$નો એક જ ઘટક
હોય, તો તેઓ~$\Domain{M}$ના એ જ ઉપગણોના ઘટકો હોય છે (કારણ
કે ગણો તેમના ઘટકો વડે નક્કી થાય છે). ઊલટું, જો~$a$ અને~$b$
જુદા હોય, તો તેઓ એકસરખા ઉપગણોના ઘટકો નથી: દાખલા તરીકે,
$a \neq b$ હોય તો $a \in \{a\}$ પરંતુ $b
\notin \{a\}$. તેથી ``$\Domain{M}$ના એ જ ઉપગણોના ઘટકો
હોવું'' એવો સંબંધ~$a$ અને~$b$ માટે ત્યારે અને માત્ર ત્યારે જ સાચો
છે જ્યારે $a = b$. $\Domain{M}$ના બધા ઉપગણો પર પરિમાણીકરણ કરી
શકાતું હોવાથી આ સંબંધ દ્વિતીય-ક્રમ તર્કશાસ્ત્રમાં વ્યક્ત કરી શકાય છે.
આમ, નીચેનું
સૂત્ર~$\Id{\Domain{M}}$ને વ્યાખ્યાયિત કરે છે:
\[
\lforall[X][(X(x) \liff X(y))]
\]
\end{ex}

\begin{prob}
બતાવો કે $\lforall[X][(X(x) \lif X(y))]$ પણ~$\Id{\Domain{M}}$
વ્યાખ્યાયિત કરે છે (ધ્યાન આપો: અહીં $\lif$ છે, $\liff$ નહિ!).
\end{prob}

\begin{ex}
જો~$R$ દ્વિસ્થાની વિધેય હોય, તો~$\Assign{R}{M}$ એ
$\Domain{M}$ પરનો દ્વિસ્થાની સંબંધ છે. થોડી ગૂંચવણ થવાની શક્યતા
હોવા છતાં, આપણે~$R$નો ઉપયોગ વિધેય~$R$ના સંકેત માટે અને
સંબંધ~$\Assign{R}{M}$ પોતે દર્શાવવા માટે પણ કરીશું. $R$નું
\emph{પરંપરિત સંવરણ}~$R^*$ એવો સંબંધ છે કે~$a$ અને~$b$ વચ્ચે તે
ત્યારે અને માત્ર ત્યારે જ હોય જ્યારે અમુક $c_1$, \dots, $c_k$
માટે $R(a,c_1)$, $R(c_1, c_2)$, \dots, $R(c_k,b)$ સાચાં હોય.
તેમાં $k = 0$નો કિસ્સો પણ છે: $R(a,b)$ હોય તો $R^*(a,b)$ પણ હોય.
તેથી $R \subseteq R^*$. હકીકતમાં~$R^*$ એવો સૌથી નાનો સંબંધ છે જે
$R$ને સમાવે છે અને પરંપરિત છે. દ્વિતીય-ક્રમ તર્કશાસ્ત્રમાં~$X$ એ~$R$
સમાવતો પરંપરિત સંબંધ છે એવું કહી શકાય છે:
\begin{multline*}
  !B_R(X) \ident \lforall[x][\lforall[y][(R(x,y) \lif X(x, y))]] \land {}\\
\lforall[x][\lforall[y][\lforall[z][((X(x,y) \land X(y,z)) \lif X(x,
      z))]]].
\end{multline*}
પહેલો સંયોજક $R \subseteq X$ કહે છે અને બીજો~$X$ પરંપરિત છે એમ કહે
છે.

$X$ એવો સૌથી નાનો સંબંધ છે એમ કહેવાનો અર્થ એ છે કે~$R$ને સમાવતા
અને પરંપરિત એવા દરેક સંબંધમાં તે પોતે સમાવિષ્ટ છે. તેથી~$R$નું
પરંપરિત સંવરણ નીચેના સૂત્ર વડે વ્યાખ્યાયિત કરી શકાય છે:
\[
R^*(X) \ident !B_R(X) \land \lforall[Y][(!B_R(Y) \lif
  \lforall[x][\lforall[y][(X(x, y) \lif Y(x,y))]])].
\]
$\Sat{M}{R^*(X)}[s]$ ત્યારે અને માત્ર ત્યારે જ જ્યારે $s(X) = R^*$.
$R$નું પરંપરિત સંવરણ પ્રથમ-ક્રમ તર્કશાસ્ત્રમાં વ્યક્ત કરી શકાતું નથી.
\end{ex}
% END OLP-0328
\endgroup


\begingroup
\makeatletter\def\ol@sectioncs{section}\makeatother

% BEGIN OLP-0329
\olfileid{sol}{syn}{siz}

\olsection{અનંત અને ગણનીય
  પ્રદેશોનું વર્ણન}
\phantomsection\label{guunit:OLP-0329}


ગણ~$M$ (ડેડેકિન્ડના અર્થમાં) અનંત ત્યારે અને માત્ર ત્યારે જ હોય જ્યારે
એવું એક-એક વિધેય~$f\colon M
\to M$ હોય જે વ્યાપ્ત ન હોય, એટલે $\ran{f} \neq M$. પ્રથમ-ક્રમ
તર્કશાસ્ત્રમાં એકસ્થાની ફલન~$f$ લઈ શકાય અને
સંરચના~$\Struct{M}$માં તેને અપાયેલું વિધેય~$\Assign{f}{M}$
એક-એક છે તથા $\ran{f} \neq
\Domain{M}$ છે એમ કહી શકાય:
\[
\lforall[x][\lforall[y][(\eq[f(x)][f(y)] \to \eq[x][y])]] \land
\lexists[y][\lforall[x][\eq/[y][f(x)]]].
\]
જો~$\Struct{M}$ આ વાક્યને સંતોષે, તો
$\Assign{f}{M}:
\Domain{M} \to \Domain{M}$ એક-એક છે અને તેથી~$\Domain{M}$
અનંત હોવું જોઈએ. જો~$\Domain{M}$ અનંત હોય અને તેથી આવું વિધેય
અસ્તિત્વમાં હોય, તો~$\Assign{f}{M}$ને તે વિધેય લઈએ; ત્યારે
$\Struct{M}$ આ વાક્યને સંતોષશે. પરંતુ આ માટે આપણા હેતુનો
અતાર્કિક સંકેત~$f$ ભાષામાં હોવો જરૂરી છે. દ્વિતીય-ક્રમ તર્કશાસ્ત્રમાં
આવું વિધેય \emph{અસ્તિત્વમાં છે} એમ સીધું કહી શકાય છે. ત્યારે~$f$ની
જરૂર રહેતી નથી અને શુદ્ધ દ્વિતીય-ક્રમ તર્કશાસ્ત્રમાં આ વાક્ય
મળે છે:
\[
\fn{Inf} \ident \lexists[u][(\lforall[x][\lforall[y][(\eq[u(x)][u(y)]
      \lif \eq[x][y])]] \land \lexists[y][\lforall[x][\eq/[y][u(x)]]])].
\]
$\Sat{M}{\fn{Inf}}$ ત્યારે અને માત્ર ત્યારે જ જ્યારે~$\Domain{M}$
અનંત હોય. પછી $\fn{Fin} \ident \lnot \fn{Inf}$ વ્યાખ્યાયિત કરી
શકાય; $\Sat{M}{\fn{Fin}}$ ત્યારે અને માત્ર ત્યારે જ જ્યારે
$\Domain{M}$ સાન્ત હોય. શુદ્ધ પ્રથમ-ક્રમ તર્કશાસ્ત્રનું એક પણ એકલું
વાક્ય પ્રદેશ અનંત છે એવું વ્યક્ત કરી શકતું નથી, જોકે
આવાં વાક્યોનો અનંત ગણ તે વ્યક્ત કરી શકે છે. શુદ્ધ પ્રથમ-ક્રમ
તર્કશાસ્ત્રનાં વાક્યોનો એવો કોઈ ગણ નથી જે કોઈ
સંરચનામાં ત્યારે અને માત્ર ત્યારે જ સંતોષાય જ્યારે તેનું
વ્યાપકક્ષેત્ર સાન્ત હોય.

\begin{prop}
$\Sat{M}{\fn{Inf}}$ ત્યારે અને માત્ર ત્યારે જ જ્યારે~$\Domain{M}$
અનંત હોય.
\end{prop}

\begin{proof}
$\Sat{M}{\fn{Inf}}$ ત્યારે અને માત્ર ત્યારે જ જ્યારે અમુક~$s$
માટે
  $\Sat{M}{\lforall[x][\lforall[y][(\eq[u(x)][u(y)] \lif \eq[x][y])]]
    \land \lexists[y][\lforall[x][\eq/[y][u(x)]]]}[s]$.
  એવું હોય તો~$s(u)$ એક-એક વિધેય છે અને અમુક $y
  \in \Domain{M}$ એ~$s(u)$ના મૂલ્યસમૂહમાં નથી. ઊલટું, જો એવું
  એક-એક $f\colon \Domain{M} \to \Domain{M}$ હોય કે
  $\ran{f} \neq \Domain{M}$, તો~$s(u) = f$ એવી ચલ-નિયુક્તિ જરૂરી
  શરત પૂરી કરે છે.
\end{proof}

અહીં અરિક્ત ગણ~$M$ને ગણનીય કહીએ છીએ જો તેના ઘટકોની
આવી પરિગણના હોય:
\[
m_0, m_1, m_2, \dots
\]
તેમાં પુનરાવર્તન નથી, પરંતુ પરિગણના સાન્ત હોઈ શકે છે.
\sourcecorrection{OLSOL-004}{સ્થિર અંગ્રેજી મૂળ અહીં મનસ્વી ગણને
પરિગણનીય કહે છે, પણ તરત પછી $m_0$ અને $z\in M$ માગે છે.
રિક્ત ગણ માટે આ શક્ય નથી. સંરચનાનાં વ્યાપકક્ષેત્રો અરિક્ત હોય છે,
અને અહીં જરૂરી અરિક્તતાની શરત સ્પષ્ટ કરી છે.}
આવી પરિગણના ત્યારે અને માત્ર ત્યારે જ મળે જ્યારે ઘટક
$z \in M$ અને વિધેય $f\colon M \to M$ એવાં હોય કે $z$, $f(z)$,
$f(f(z))$, \dots, ના જુદા જુદા ઘટકો જ~$M$ના બધા ઘટકો હોય.
કારણ કે જો પરિગણના હોય, તો $z = m_0$ અને $f(m_k) = m_{k+1}$
(અથવા~$m_k$ પરિગણનાનો છેલ્લો ઘટક હોય તો
$f(m_k) = m_k$) જરૂરી ઘટક અને વિધેય આપે છે. બીજી બાજુ,
આવાં~$z$ અને~$f$ હોય તો $z$, $f(z)$, $f(f(z))$, \dots, માંથી
પુનરાવર્તિત ઘટકો કાઢતાં~$M$ની પરિગણના મળે છે અને~$M$
ગણનીય છે. આ~$z$ અને~$f$નું અસ્તિત્વ દ્વિતીય-ક્રમ
તર્કશાસ્ત્રમાં વ્યક્ત કરીને એવું વાક્ય બનાવી શકાય છે જે
સંરચનામાં ત્યારે અને માત્ર ત્યારે જ સાચું હોય જ્યારે તે
સંરચના ગણનીય હોય:
\[
\fn{Count} \ident
\lexists[z][\lexists[u][\lforall[X][((X(z) \land
      \lforall[x][(X(x) \lif X(u(x)))]) \lif \lforall[x][X(x)])]]]
\]
\sourcecorrection{OLSOL-005}{સ્થિર અંગ્રેજી મૂળમાં પરિગણના
પુનરાવર્તન વિનાની કહેવાઈ છે, પણ સાન્ત પરિગણનાના છેલ્લા ઘટક પર
વિધેય સ્થિર થાય ત્યારે તેની આગળની કક્ષા એ ઘટકને વારંવાર આપે છે.
અહીં કક્ષાના જુદા ઘટકો પરિગણના બનાવે છે એમ સ્પષ્ટ કર્યું છે;
નીચેની સાબિતીમાં પણ છેલ્લા ઘટકનું વિધેય-મૂલ્ય સ્પષ્ટ કરાયું છે.}

\begin{prop}
$\Sat{M}{\fn{Count}}$ ત્યારે અને માત્ર ત્યારે જ જ્યારે~$\Domain{M}$
ગણનીય હોય.
\end{prop}

\begin{proof}
માનો કે~$\Domain{M}$ ગણનીય છે અને $m_0$, $m_1$, \dots,
તેની પરિગણના છે. પુનરાવર્તનો દૂર કરીને કોઈ~$m_k$ બે વાર ન આવે તેની
ખાતરી કરી શકાય છે. $f(m_k) = m_{k+1}$ વ્યાખ્યાયિત કરો (સાન્ત
પરિગણનાના છેલ્લા ઘટકને પોતાને જ મોકલો) અને $s(z) = m_0$ તથા
$s(u) = f$ લો. આપણે બતાવીએ કે
\[
\Sat{M}{\lforall[X][((X(z) \land \lforall[x][(X(x) \lif X(u(x)))])
    \lif \lforall[x][X(x)])]}[s]
\]
માનો કે $M \subseteq \Domain{M}$ મનસ્વી છે. વધુમાં માનો કે
$\Sat{M}{(X(z) \land \lforall[x][(X(x) \lif
X(u(x)))])}[\Subst{s}{M}{X}]$. ત્યારે
$\Subst{s}{M}{X}(z) \in M$ અને જ્યારે પણ $x \in M$, ત્યારે
$(\Subst{s}{M}{X}(u))(x) \in M$ પણ. બીજા શબ્દોમાં,
$\varAssign{\Subst{s}{M}{X}}{s}{X}$ હોવાથી, $m_0 \in M$ અને
$x \in M$ હોય તો $f(x) \in M$. તેથી $m_0 \in M$,
$m_1 = f(m_0) \in M$, $m_2 = f(f(m_0)) \in M$ વગેરે. આમ,
$M = \Domain{M}$ અને તેથી
$\Sat{M}{\lforall[x][X(x)]}[\Subst{s}{M}{X}]$. $M \subseteq
\Domain{M}$ મનસ્વી હોવાથી સાબિતી પૂરી થાય છે:
$\Sat{M}{\fn{Count}}$.

હવે માનો કે $\Sat{M}{\fn{Count}}$, એટલે અમુક~$s$ માટે
\[
\Sat{M}{\lforall[X][((X(z) \land \lforall[x][(X(x) \lif X(u(x)))])
    \lif \lforall[x][X(x)])]}[s]
\]
$m = s(z)$ અને $f = s(u)$ લો તથા
$M = \{m,
f(m), f(f(m)), \dots\}$ વિચારો. પુનરાવર્તિત ઘટકો કાઢીએ તો
આ~$M$ ગણનીય છે. ત્યારે ધારણા પરથી
\[
\Sat{M}{(X(z) \land \lforall[x][(X(x) \lif X(u(x)))] )
    \lif \lforall[x][X(x)]}[\Subst{s}{M}{X}]
\]
મળે છે. વળી $M \ni m =
\Subst{s}{M}{X}(z)$ હોવાથી $\Sat{M}{X(z)}[\Subst{s}{M}{X}]$,
અને $x \in M$ હોય ત્યારે $f(x) \in
M$ હોવાથી $\Sat{M}{\lforall[x][(X(x) \lif
X(u(x)))]}[\Subst{s}{M}{X}]$ પણ મળે છે. આમ, પૂર્વપક્ષ અને
શરતી વિધાન બંને સંતોષાય છે, તેથી ઉત્તરપક્ષ પણ સંતોષાવો જોઈએ:
$\Sat{M}{\lforall[x][X(x)]}[\Subst{s}{M}{X}]$. પરંતુ તેનો અર્થ
$M = \Domain{M}$; અને~$M$ વ્યાખ્યા મુજબ પરિગણનીય હોવાથી
$\Domain{M}$ પણ ગણનીય છે.
\end{proof}

\begin{prob}
વાક્ય~$\fn{Inf} \land \fn{Count}$ બધા અને માત્ર
ગણનીય વ્યાપકક્ષેત્રોમાં સાચું છે. $\fn{Count}$ની વ્યાખ્યા
એવી બદલો કે નવું વાક્ય વ્યાપકક્ષેત્ર ગણનીય છે એવું
સીધું વ્યક્ત કરે, અને તે કરે છે એમ સાબિત કરો.
\end{prob}
% END OLP-0329
\endgroup


\OLEndChapterHook


\begin{editorial}
આગળનો સંબંધિત અંશ મૂળના સક્રિય આયાતક્રમમાં નથી. તેને અહીં સ્પષ્ટ વૈકલ્પિક સંદર્ભમાં જાળવ્યો છે.
\end{editorial}
\begingroup
\makeatletter\def\ol@sectioncs{section}\makeatother

% BEGIN OLP-0716
\olfileid{sol}{syn}{lan}

\olsection{દ્વિતીય-ક્રમ તર્કશાસ્ત્રની ભાષા}
\phantomsection\label{guunit:OLP-0716}


પ્રથમ-ક્રમ તર્કશાસ્ત્રની જેમ, દ્વિતીય-ક્રમ તર્કશાસ્ત્રના
વ્યંજકો મૂળભૂત શબ્દભંડોળમાંથી રચાય છે. તેમાં
\emph{ચલો}, \emph{અચળો},
\emph{વિધેયો} અને ક્યારેક \emph{ફલનો}
હોય છે. તેમનાથી, તાર્કિક સંયોજકો, પરિમાણકો અને કૌંસ તથા
અલ્પવિરામ જેવાં વિરામચિહ્નો સાથે, \emph{પદો} અને
\emph{સૂત્રો} રચાય છે.
ફરક એટલો છે કે વસ્તુઓ માટેના ચલો ઉપરાંત દ્વિતીય-ક્રમ
તર્કશાસ્ત્રમાં સંબંધો અને ફલનો માટેના ચલો પણ છે,
અને તેમના પર પરિમાણીકરણની છૂટ છે.
એટલે દ્વિતીય-ક્રમ તર્કશાસ્ત્રનાં તાર્કિક ચિહ્નો
પ્રથમ-ક્રમનાં ચિહ્નો ઉપરાંત આ છે:

\begin{enumerate}
\item દરેક અરીટી~$n$ માટે દ્વિતીય-ક્રમ સંબંધ
ચલોનો ગણનીય ગણ:
$\Obj V_0^n$, $\Obj V_1^n$, $\Obj V_2^n$, \dots
\item દ્વિતીય-ક્રમ ફલન ચલોનો ગણનીય ગણ:
$\Obj u_0^n$, $\Obj u_1^n$, $\Obj u_2^n$, \dots
\end{enumerate}

પ્રથમ-ક્રમ તર્કશાસ્ત્રમાં સામાન્ય રીતે તાદાત્મ્ય~$\eq$
સમાવાય છે. તેમાં ભાષા~$\Lang{L}$નાં અતાર્કિક ચિહ્નો
રસપ્રદ વાતો વ્યક્ત કરવા માટે અગત્યનાં છે.
અતાર્કિક ચિહ્નો વિનાનાં વાક્યો પણ છે,
પરંતુ ફક્ત $\eq$થી રસપ્રદ વાતો કહેવી મુશ્કેલ છે.
દ્વિતીય-ક્રમ તર્કશાસ્ત્રમાં સંબંધ અને ફલનના ચલોનો અમર્યાદિત
પુરવઠો હોવાથી, વિશેષ અતાર્કિક ચિહ્નોને અનુરૂપ ચલો વડે
પ્રથમ-ક્રમ ભાષાની વાતો રજૂ કરી શકીએ;
અર્થ જાળવવા આ ચલોને તે ચિહ્નોનાં અર્થઘટનનાં મૂલ્યો આપવા પડે છે.
\sourcecorrection{OLMD-346}{મૂળના છેલ્લાં વાક્યમાં અતાર્કિક ચિહ્નો વિના પ્રથમ-ક્રમ ભાષાની કોઈ પણ વાત કહી શકાય તેવો અમર્યાદિત દાવો છે. સંબંધ અને ફલનના મુક્ત ચલો તે ચિહ્નોની જગ્યાએ મૂકી, તેમના અર્થઘટનને અનુરૂપ નિયુક્તિ આપતાં સંતોષ જાળવાય છે; અચળ માટે વસ્તુ-ચલ વાપરી શકાય છે. આ ચલો પર આપમેળે પરિમાણીકરણ કરી દેવાથી મૂળનું ચોક્કસ અર્થઘટન જળવાતું નથી. તેથી મુક્ત ચલો અને નિયુક્તિની જરૂરી શરત સ્પષ્ટ કરી છે.}
% END OLP-0716
\endgroup
% END OLP-0323
\endgroup


\begingroup
\makeatletter\def\ol@sectioncs{section}\makeatother

% BEGIN OLP-0330
\olchapter{sol}{met}{દ્વિતીય-ક્રમ તર્કશાસ્ત્રનો અધિસિદ્ધાંત}
\olchapterid{sol}{met}
\phantomsection\label{guunit:OLP-0330}


\begingroup
\makeatletter\def\ol@sectioncs{section}\makeatother

% BEGIN OLP-0331
\olfileid{sol}{met}{int}

\olsection{પ્રસ્તાવના}
\phantomsection\label{guunit:OLP-0331}


પ્રથમ-ક્રમ તર્કશાસ્ત્રમાં કેટલાક સારા ગુણધર્મો છે. આપણે જાણીએ છીએ કે
તે નિર્ણેય નથી, પરંતુ ઓછામાં ઓછું સ્વયંસિદ્ધીકરણીય છે. એટલે કે
પ્રથમ-ક્રમ તર્કશાસ્ત્ર માટે યથાર્થ અને પૂર્ણ સાબિતી-તંત્રો છે: તેઓ
એવો નિષ્પન્નક્ષમતા સંબંધ~$\Proves$ આપે છે કે વાક્યોના કોઈપણ
ગણ~$\Gamma$ અને વાક્ય~$!A$ માટે $\Gamma \Entails !A$ ત્યારે અને માત્ર
ત્યારે જ જ્યારે $\Gamma \Proves !A$. ખાસ કરીને તેનો અર્થ એ થાય છે કે
પ્રથમ-ક્રમ તર્કશાસ્ત્રનાં પ્રમાણભૂત વાક્યો સંગણકીય રીતે પરિગણનીય છે. એવું સંગણનીય વિધેય~$f\colon \Nat \to
\Sent[L]$ છે કે~$f$નાં મૂલ્યો~$\Lang{L}$નાં બધાં અને માત્ર પ્રમાણભૂત
વાક્યો છે. આવું બને છે કારણ કે નિષ્પત્તિઓની પરિગણના
કરી શકાય છે અને એક વાક્ય નિષ્પન્ન કરતી નિષ્પત્તિઓને પછી
તે વાક્ય સાથે જોડી શકાય છે. દ્વિતીય-ક્રમ તર્કશાસ્ત્ર
પ્રથમ-ક્રમ તર્કશાસ્ત્ર કરતાં વધુ અભિવ્યક્તિશીલ છે; તેથી તેનાં
પ્રમાણભૂત વાક્યોનું નિરૂપણ સામાન્ય રીતે વધુ જટિલ છે. હકીકતમાં આપણે
બતાવીશું કે દ્વિતીય-ક્રમ તર્કશાસ્ત્ર માત્ર અનિર્ણેય જ નથી, પણ તેના
પ્રમાણભૂત વાક્યો સંગણકીય રીતે પરિગણનીય પણ નથી. તેથી દ્વિતીય-ક્રમ તર્કશાસ્ત્ર માટે યથાર્થ અને
પૂર્ણ સાબિતી-તંત્ર હોઈ શકે નહિ (જોકે યથાર્થ પરંતુ અપૂર્ણ સાબિતી-તંત્રો
ઉપલબ્ધ છે અને તે સંશોધનના મહત્વના વિષયો પણ છે).

પ્રથમ-ક્રમ તર્કશાસ્ત્રના બીજા બે ગુણધર્મો પણ છે: તે સઘન છે (વાક્યોના
ગણ~$\Gamma$નો દરેક સાન્ત ઉપગણ સંતોષ્ય હોય તો~$\Gamma$ પોતે સંતોષ્ય
છે) અને તેના માટે લેવેનહાઇમ--સ્કોલેમ પ્રમેય સાચું છે (જો~$\Gamma$નું
અનંત નિદર્શ હોય તો તેનું ગણનીય નિદર્શ પણ હોય છે).
દ્વિતીય-ક્રમ તર્કશાસ્ત્રમાં આ બંને પરિણામો નિષ્ફળ જાય છે. અહીં પણ
કારણ એ જ છે કે દ્વિતીય-ક્રમ તર્કશાસ્ત્ર પ્રદેશોના કદ વિશે એવી
હકીકતો વ્યક્ત કરી શકે છે જે પ્રથમ-ક્રમ તર્કશાસ્ત્ર કરી શકતું નથી.
% END OLP-0331
\endgroup


\begingroup
\makeatletter\def\ol@sectioncs{section}\makeatother

% BEGIN OLP-0332
\olfileid{sol}{met}{spa}

\olsection{દ્વિતીય-ક્રમ અંકગણિત}
\phantomsection\label{guunit:OLP-0332}


યાદ કરો કે પેયાનો અંકગણિતના સિદ્ધાંત~$\Th{PA}$માં~$\Th{Q}$નાં આઠ
સ્વયંસિદ્ધિઓ,
\begin{align*}
& \lforall[x][\eq/[x'][\Obj 0]]\\
& \lforall[x][\lforall[y][(\eq[x'][y'] \lif \eq[x][y])]]\\
& \lforall[x][(\eq[x][\Obj 0] \lor \lexists[y][\eq[x][y']])]\\
& \lforall[x][\eq[(x + \Obj 0)][x]]\\
& \lforall[x][\lforall[y][\eq[(x + y')][(x + y)']]]\\
& \lforall[x][\eq[(x \times \Obj 0)][\Obj 0]]\\
& \lforall[x][\lforall[y][\eq[(x \times y')][((x \times y) + x)]]]\\
& \lforall[x][\lforall[y][(x < y \liff \lexists[z][\eq[(z' + x)][y]])]]\\
\intertext{અને આ સ્વરૂપનાં બધાં વાક્યો સામેલ છે:}
& (!A(\Obj 0) \land \lforall[x][(!A(x) \lif !A(x'))]) \lif \lforall[x][!A(x)].
\end{align*}
છેલ્લું ``પ્રરૂપ'' છે, એટલે અંકગણિતની ભાષાના દરેક
સૂત્ર~$!A(x)$ માટે એક, એમ અનંત સંખ્યામાં વાક્યો
ઉત્પન્ન કરતું માળખું. આ પ્રરૂપને (પ્રથમ-ક્રમ) \emph{આગમનનું
સ્વયંસિદ્ધિ-પ્રરૂપ} કહીએ છીએ.
\emph{દ્વિતીય-ક્રમ} પેયાનો અંકગણિત~$\Th{PA^2}$માં આગમનને એક જ
વાક્યથી વ્યક્ત કરી શકાય છે. $\Th{PA^2}$ ઉપરની પ્રથમ આઠ સ્વયંસિદ્ધિઓ
અને (દ્વિતીય-ક્રમ) \emph{આગમનના સ્વયંસિદ્ધિ}થી બને છે:
\[
\lforall[X][((X(\Obj 0) \land \lforall[x][(X(x) \lif X(x'))]) \lif \lforall[x][X(x)])].
\]
તે કહે છે: પ્રદેશનો ઉપગણ~$X$ જો~$\Assign{\Obj{0}}{M}$ને સમાવે
અને કોઈપણ~$x \in \Domain{M}$ સાથે~$\Assign{\prime}{M}(x)$ને પણ
સમાવે (એટલે કે તે ``અનુગામી હેઠળ બંધ'' હોય), તો તે આખું
પ્રદેશ સમાવે (એટલે $X =
\Domain{M})$.

આગમનનું સ્વયંસિદ્ધિ ખાતરી આપે છે કે તેને સંતોષતી કોઈપણ સંરચનાના
$\Domain{M}$માં સ્વયંસિદ્ધિઓથી હોવા જરૂરી હોય એવા ઘટકો જ છે,
એટલે $n \in \Nat$ માટે~$\num{n}$નાં મૂલ્યો. $\Th{PA^2}$ના
નિદર્શમાં કોઈ અપ્રમાણભૂત સંખ્યાઓ નથી.

\begin{thm}
\ollabel{thm:sol-pa-standard}
જો $\Sat{M}{\Th{PA^2}}$ હોય, તો $\Domain{M} =
\Setabs{\Value{\num{n}}{M}}{n \in \Nat}$.
\end{thm}

\begin{proof}
$N = \Setabs{\Value{\num{n}}{M}}{n \in \Nat}$ લો અને
$\Sat{M}{\Th{PA^2}}$ માનો. દરેક $n \in \Nat$ માટે
$\Value{\num{n}}{M} \in \Domain{M}$ છે જ; તેથી
$N \subseteq \Domain{M}$.

હવે બીજી દિશાનો સમાવેશ બતાવીએ. $s(X) = N$ એવી ચલ-નિયુક્તિ~$s$
વિચારો. ધારણા મુજબ,
\begin{align*}
\Sat{M}{\lforall[X][((X(\Obj 0) \land \lforall[x][(X(x)
      \lif X(x'))]) \lif \lforall[x][X(x)])]}, & \text{ તેથી}\\
\Sat{M}{(X(\Obj 0) \land \lforall[x][(X(x) \lif X(x'))]) \lif
  \lforall[x][X(x)]}[s]. &
\end{align*}
આ શરતી વિધાનનો પૂર્વપક્ષ વિચારો. $\Value{\Obj{0}}{M} \in
N$ છે; તેથી $\Sat{M}{X(\Obj{0})}[s]$. બીજો સંયોજક,
$\lforall[x][(X(x) \lif X(x'))]$, પણ સંતોષાય છે. કારણ કે
$x \in N$ હોય એમ માનો. $N$ની વ્યાખ્યા મુજબ કોઈ~$n$ માટે
$x = \Value{\num{n}}{M}$. તેથી
$\Assign{\prime}{M}(x) = \Value{\num{n+1}}{M} \in N$. આમ,
$\Assign{\prime}{M}(x) \in N$.

આથી $\Sat{M}{X(\Obj 0) \land \lforall[x][(X(x) \lif X(x'))]}[s]$.
પરિણામે $\Sat{M}{\lforall[x][X(x)]}[s]$. પરંતુ તેનો અર્થ એ છે કે
દરેક $x \in \Domain{M}$ માટે $x \in s(X) = N$. તેથી
$\Domain{M} \subseteq N$.
\end{proof}

\begin{cor}
\ollabel{cor:sol-pa-categorical}
$\Th{PA^2}$નાં કોઈપણ બે નિદર્શો એકરૂપ છે.
\end{cor}

\begin{proof}
\olref{thm:sol-pa-standard} મુજબ $\Th{PA^2}$ના કોઈપણ નિદર્શનું
વ્યાપકક્ષેત્ર $\Value{\num{n}}{M}$નાં મૂલ્યોથી પૂરું થાય છે. આવું
દરેક નિદર્શ~$\Th{Q}$નું નિદર્શ પણ છે.
\olref[mod][mar][stm]{prop:thq-standard} મુજબ આવું દરેક નિદર્શ
પ્રમાણભૂત છે, એટલે~$\Struct{N}$ને એકરૂપ છે.
\end{proof}

ઉપર આપણે~$\Th{PA^2}$ને પ્રથમ આઠ અંકગણિતીય સ્વયંસિદ્ધિઓ તથા
દ્વિતીય-ક્રમ આગમનનું સ્વયંસિદ્ધિ ધરાવતા સિદ્ધાંત તરીકે વ્યાખ્યાયિત
કર્યો. હકીકતમાં દ્વિતીય-ક્રમ તર્કશાસ્ત્રની અભિવ્યક્તિની શક્તિને
કારણે દ્વિતીય-ક્રમ પેયાનો અંકગણિત માટે અંકગણિતીય સ્વયંસિદ્ધિઓમાંથી
માત્ર \emph{પ્રથમ બે} અને આગમન પૂરતાં છે.

\begin{prop}\ollabel{prop:sol-pa-definable}
$\Th{PA^{2\dagger}}$ એ પ્રથમ બે અંકગણિતીય સ્વયંસિદ્ધિઓ (અનુગામી
વિશેની સ્વયંસિદ્ધિઓ) અને દ્વિતીય-ક્રમ આગમનનું સ્વયંસિદ્ધિ ધરાવતો
દ્વિતીય-ક્રમ સિદ્ધાંત હોય. ત્યારે $\le$, $+$ અને $\times$
$\Th{PA^{2\dagger}}$ના અનુગામી-માત્ર પાયામાં વ્યાખ્યાયિત કરી શકાય છે.
\end{prop}
\sourcecorrection{OLSOL-006}{સ્થિર અંગ્રેજી મૂળમાં આ ક્રિયાઓ
``સિદ્ધાંતમાં વ્યાખ્યેય'' છે એમ કહેવાયું છે. અર્થ એવો છે કે
અનુગામી અને શૂન્ય ધરાવતા પાયામાં દ્વિતીય-ક્રમ સૂત્રો વડે પ્રમાણભૂત
ક્રમ, સરવાળો અને ગુણાકારના સંબંધો નિર્ધારિત કરી શકાય છે. પ્રથમ બે
સ્વયંસિદ્ધિઓ મનસ્વી રીતે પહેલેથી અર્થઘટિત $+$ કે $\times$ને
નિયંત્રિત કરતી નથી.}

\begin{proof}
$\le$ વ્યાખ્યાયિત છે એમ બતાવવા એવું સૂત્ર~$!A_\le(x, y)$ શોધવું
પડે કે $\Sat{N}{!A_\le(\num n, \num m)}$ ત્યારે અને માત્ર ત્યારે જ
જ્યારે $n \le m$. આ સૂત્ર વિચારો:
\[
!B(x, Y) \ident Y(x) \land \lforall[y][(Y(y) \lif Y(y'))]
\]
સ્પષ્ટ છે કે $!B(\num n, Y)$ કોઈ ગણ $Y \subseteq \Nat$થી ત્યારે
અને માત્ર ત્યારે જ સંતોષાય જ્યારે $\Setabs{m}{n \le m} \subseteq Y$.
તેથી $!A_\le(x, y) \ident
\lforall[Y][(!B(x, Y) \lif Y(y))]$ લઈ શકાય.

સરવાળો વ્યાખ્યાયિત છે એ જોવા નોંધો કે $k+l=m$ ત્યારે અને માત્ર ત્યારે
જ જ્યારે એવું વિધેય~$u$ હોય કે $u(0)=k$, બધા~$n$ માટે
$u(n')=u(n)'$ અને $m = u(l)$. આ સમકક્ષતા વાપરીને
$\Th{PA^{2\dagger}}$માં સરવાળો નીચેના સૂત્ર વડે વ્યાખ્યાયિત કરી
શકાય:
\[
!A_+(x,y,z) \ident \lexists[u][(u(\Obj 0)=x \land \lforall[w][u(w')=u
 (w)'] \land u(y)=z)]
\]
હવે $\Sat{N}{!A_+(\num k, \num l, \num m)}$ ત્યારે અને માત્ર
ત્યારે જ જ્યારે $k+l=m$ એ સ્પષ્ટ હોવું જોઈએ.
\sourcecorrection{OLSOL-007}{સ્થિર અંગ્રેજી મૂળ સરવાળાના સૂત્રમાં
$\lforall[w]$ લખે છે, પરંતુ તેની અંદર પુનરાવર્તન માટે $u(x')=u(x)'$
લખે છે. તેથી $w$ વપરાતો નથી અને જરૂરી સર્વવાચી પુનરાવર્તન મળતું
નથી. અહીં $u(w')=u(w)'$ કર્યું છે.}
\end{proof}

\begin{prob}
\olref[sol][met][spa]{prop:sol-pa-definable}ની સાબિતી પૂરી કરો.
\end{prob}
% END OLP-0332
\endgroup


\begingroup
\makeatletter\def\ol@sectioncs{section}\makeatother

% BEGIN OLP-0333
\olfileid{sol}{met}{nax}

\olsection{દ્વિતીય-ક્રમ તર્કશાસ્ત્ર સ્વયંસિદ્ધીકરણીય નથી}
\phantomsection\label{guunit:OLP-0333}



\begin{thm}
\ollabel{thm:sol-undecidable}
દ્વિતીય-ક્રમ તર્કશાસ્ત્ર અનિર્ણેય છે.
\end{thm}

\begin{proof}
પ્રથમ-ક્રમ વાક્ય પ્રથમ-ક્રમ તર્કશાસ્ત્રમાં પ્રમાણભૂત ત્યારે
અને માત્ર ત્યારે જ હોય જ્યારે તે દ્વિતીય-ક્રમ તર્કશાસ્ત્રમાં
પ્રમાણભૂત હોય; વળી પ્રથમ-ક્રમ તર્કશાસ્ત્ર અનિર્ણેય છે.
\end{proof}

\begin{thm}
\ollabel{cor:sol-not-axiomatizable}
દ્વિતીય-ક્રમ તર્કશાસ્ત્ર માટે અસરકારક રીતે આપેલું એવું કોઈ યથાર્થ અને
પૂર્ણ નિષ્પત્તિ તંત્ર નથી.
\end{thm}
\sourcecorrection{OLSOL-008}{સ્થિર અંગ્રેજી મૂળમાં અહીં કોઈપણ યથાર્થ
અને પૂર્ણ નિષ્પત્તિ-તંત્ર નથી એમ નિઃશરત કહેવાયું છે. અસરકારક રીતે
આપેલું તંત્ર હોવાની શરત જરૂરી છે: અસરકારકતાની શરત વિના બધાં પ્રમાણભૂત
વાક્યોને જ સ્વયંસિદ્ધિઓ તરીકે લઈ શકાય. નીચેની સાબિતી સંગણનીય
પરિગણનાનો વિરોધાભાસ ઉત્પન્ન કરે છે, તેથી આ શરત સ્પષ્ટ કરી છે.}

\begin{proof}
$!A$ અંકગણિતની ભાષાનું વાક્ય હોય. $\Sat{N}{!A}$ ત્યારે
અને માત્ર ત્યારે જ જ્યારે $\Th{PA^2} \Entails !A$. $!P$ એ
$\Th{PA^2}$નાં નવ સ્વયંસિદ્ધિઓનો સંયોજક હોય. $\Th{PA^2} \Entails !A$
ત્યારે અને માત્ર ત્યારે જ જ્યારે $\Entails !P \lif
!A$, એટલે દરેક સંરચના માટે
$\Sat{M}{!P \lif !A}$. 
\sourcecorrection{OLSOL-009}{સ્થિર અંગ્રેજી મૂળના આ સંતોષ-સૂત્રમાં
બીજી \texttt{\textbackslash Sat} દલીલનું બંધ કૌંસ ગણિત-મર્યાદાની
બહાર મૂકાયું છે, તેથી મૂળ TeX અસંતુલિત છે. અહીં
$\Sat{M}{!P \lif !A}$નું બંધ કૌંસ સૂત્રની અંદર મૂક્યું અને
પહેલાંના પ્રામાણ્ય-દાવાને અનુરૂપ દરેક સંરચનાનો વ્યાપ સ્પષ્ટ કર્યો છે.}
હવે~$\Obj{0}$ની જગ્યાએ~$z$, $\prime$ની જગ્યાએ એકસ્થાની
વિધેય-ચલ~$u$, $+$ અને~$\times$ની જગ્યાએ અનુક્રમે દ્વિસ્થાની
વિધેય-ચલો~$u'$ અને~$u''$, તથા $<$ની જગ્યાએ દ્વિસ્થાની સંબંધ-ચલ~$L$
મૂકીને, પછી આ બધા ચલો પર સર્વવાચી પરિમાણીકરણ કરતાં મળતું શુદ્ધ
દ્વિતીય-ક્રમ વાક્ય
$\lforall[z][\lforall[u][\lforall[u'][\lforall[u''][\lforall[L][(!P' \lif !A')]]]]]$
વિચારો. આ વાક્ય પ્રમાણભૂત ત્યારે અને માત્ર ત્યારે જ જ્યારે મૂળ વાક્ય
પ્રમાણભૂત હોય; તે ત્યારે અને માત્ર ત્યારે જ જ્યારે
$\Th{PA^2} \Entails !A$, અને તે ત્યારે અને માત્ર ત્યારે જ જ્યારે
$\Sat{N}{!A}$. તેથી દ્વિતીય-ક્રમ તર્કશાસ્ત્ર માટે યથાર્થ અને પૂર્ણ
અસરકારક સાબિતી-તંત્ર હોત, તો તેની સાબિતીઓની પરિગણનામાંથી આ ખાસ
આકારનાં વાક્યો પસંદ કરીને~$\Struct{N}$માં સાચાં વાક્યોની
સંગણનીય પરિગણના~$f\colon \Nat \to \Sent[L_A]$ મેળવી શકાત. આ
વિધેય~$\Th{Q}$માં કોઈ પ્રથમ-ક્રમ સૂત્ર~$!B_f(x, y)$ વડે નિરૂપણક્ષમ
હોત. ત્યારે સૂત્ર~$\lexists[x][!B_f(x, y)]$ એ
$\Struct{N}$માં સાચાં પ્રથમ-ક્રમ વાક્યોની ગોડેલ-સંખ્યાઓનો
ગણ વ્યાખ્યાયિત કરત, જે ટાર્સ્કીના પ્રમેયનો વિરોધ કરે છે.
\sourcecorrection{OLSOL-010}{સ્થિર અંગ્રેજી મૂળ વાક્યોની સંગણનીય
પરિગણનાથી અંકગણિતમાં વ્યાખ્યેય ગણ સુધીના પગલામાં વાક્યોની સંખ્યાકીય
સંકેતીકરણ-પ્રક્રિયા સ્પષ્ટ કરતું નથી. અહીં $y$ને સાચા વાક્યની
ગોડેલ-સંખ્યા તરીકે ઓળખાવી છે; વાક્ય પોતે અંકગણિતના ચલનું મૂલ્ય નથી.}
\end{proof}
% END OLP-0333
\endgroup


\begingroup
\makeatletter\def\ol@sectioncs{section}\makeatother

% BEGIN OLP-0334
\olfileid{sol}{met}{com}

\olsection{દ્વિતીય-ક્રમ તર્કશાસ્ત્ર સઘન નથી}
\phantomsection\label{guunit:OLP-0334}


\begin{explain}
વાક્યોનો ગણ~$\Gamma$ \emph{સાન્ત રીતે સંતોષ્ય} ત્યારે કહેવાય જ્યારે
તેનો દરેક સાન્ત ઉપગણ સંતોષ્ય હોય. પ્રથમ-ક્રમ તર્કશાસ્ત્રમાં ગુણધર્મ છે
કે વાક્યોનો ગણ~$\Gamma$ સાન્ત રીતે સંતોષ્ય હોય તો તે પોતે
સંતોષ્ય છે. આ ગુણધર્મને \emph{સઘનતા} કહે છે. ફલિતતા દ્વારા તેનું
સમકક્ષ સ્વરૂપ છે: $\Gamma \Entails !A$ હોય તો કોઈ સાન્ત ઉપગણ
$\Gamma_0 \subseteq \Gamma$ માટે પણ $\Gamma_0 \Entails
!A$. આ સ્વરૂપમાં તે પૂર્ણતા પ્રમેયનું તત્કાલ અનુસિદ્ધાંત છે: જો
$\Gamma \Entails !A$, તો પૂર્ણતા મુજબ $\Gamma \Proves !A$. પરંતુ
નિષ્પત્તિ,~$\Gamma$નાં માત્ર સાન્ત સંખ્યામાં વાક્યોનો
ઉપયોગ કરી શકે છે.

દ્વિતીય-ક્રમ તર્કશાસ્ત્રમાં સઘનતા સાચી નથી. દ્વિતીય-ક્રમ
વાક્યોના એવા ગણો છે જે સાન્ત રીતે સંતોષ્ય છે પરંતુ પોતે
સંતોષ્ય નથી; અને એવા ગણો પણ છે જે અમુક~$!A$ને ફલિત કરે છે,
પરંતુ તેમનો કોઈ સાન્ત ઉપગણ~$!A$ને ફલિત કરતો નથી.
\end{explain}


\begin{thm}
\ollabel{thm:sol-not-compact}
દ્વિતીય-ક્રમ તર્કશાસ્ત્ર સઘન નથી.
\end{thm}
\sourcecorrection{OLSOL-011}{સ્થિર અંગ્રેજી મૂળમાં આ પ્રમેયનું
\texttt{thm:sol-undecidable} લેબલ અગાઉના અનિર્ણેયતા પ્રમેયનું લેબલ
જ ફરી વાપરે છે. બે અલગ પરિણામોને એક જ લેબલ ન રહે તે માટે અહીં
\texttt{thm:sol-not-compact} આપ્યું છે.}

\begin{proof}
યાદ કરો કે
\[
\fn{Inf} \ident \lexists[u][(\lforall[x][\lforall[y][(\eq[u(x)][u(y)]
      \lif \eq[x][y])]] \land \lexists[y][\lforall[x][\eq/[y][u(x)]]])]
\]
એવું વાક્ય છે જે સંરચનામાં ત્યારે અને માત્ર ત્યારે જ
સંતોષાય જ્યારે તેનું વ્યાપકક્ષેત્ર અનંત હોય. $!A^{\ge n}$ એવું
વાક્ય હોય જે વ્યાપકક્ષેત્રમાં ઓછામાં ઓછા $n$ ઘટકો છે એમ કહે;
દાખલા તરીકે,
\[
!A^{\ge n} \ident \lexists[x_1][\dots\lexists[x_n][(\eq/[x_1][x_2]
    \land \eq/[x_1][x_3] \land \dots \land \eq/[x_{n-1}][x_n])]].
\]
આ વાક્યોનો ગણ વિચારો:
\[
\Gamma = \{\lnot \fn{Inf}, !A^{\ge 1}, !A^{\ge 2}, !A^{\ge 3}, \dots\}.
\]
તે સાન્ત રીતે સંતોષ્ય છે: મનસ્વી સાન્ત ઉપગણ~$\Gamma_0
\subseteq \Gamma$ માટે ઓછામાં ઓછો એક $k$ એવો પસંદ કરી શકાય કે
$!A^{\ge k} \in \Gamma$ અને $n>k$ માટે પસંદ કરેલા સાન્ત ઉપગણમાં
કોઈ $!A^{\ge n}
\in \Gamma_0$ ન હોય. જો $\Domain{M}$માં $k$ ઘટકો હોય,
તો $\Sat{M}{\Gamma_0}$. પરંતુ~$\Gamma$ સંતોષ્ય નથી: જો
$\Sat{M}{\lnot \fn{Inf}}$, તો~$\Domain{M}$ સાન્ત હોવું જોઈએ;
માનો તેનું કદ~$k$ છે. ત્યારે $\Sat/{M}{!A^{\ge k+1}}$.
\sourcecorrection{OLSOL-012}{સ્થિર અંગ્રેજી મૂળ સાન્ત ઉપગણની
દલીલમાં $n>k$ માટે $!A^{\ge n}$ આખા $\Gamma$માં નથી એમ કહે છે,
પરંતુ એ બધા વાક્યો $\Gamma$માં છે. અહીં તે શરત પસંદ કરેલા સાન્ત
ઉપગણ~$\Gamma_0$ માટે મૂકી છે. તેમાં કદની એક પણ શરત ન હોય તો
ઓછામાં ઓછું એક ઘટક ધરાવતું નિદર્શ લેવા $k$ને એક લઈ શકાય છે.}
\end{proof}

\begin{prob}
એવો ગણ~$\Gamma$ અને વાક્ય~$!A$નું ઉદાહરણ આપો કે
$\Gamma \Entails !A$, પરંતુ દરેક સાન્ત ઉપગણ~$\Gamma_0 \subseteq
\Gamma$ માટે $\Gamma_0 \Entails/ !A$.
\end{prob}
% END OLP-0334
\endgroup


\begingroup
\makeatletter\def\ol@sectioncs{section}\makeatother

% BEGIN OLP-0335
\olfileid{sol}{met}{lst}

\olsection{દ્વિતીય-ક્રમ તર્કશાસ્ત્રમાં લેવેનહાઇમ--સ્કોલેમ પ્રમેય નિષ્ફળ જાય છે}
\phantomsection\label{guunit:OLP-0335}


\begin{explain}
(અધોગામી) લેવેનહાઇમ--સ્કોલેમ પ્રમેય કહે છે કે વાક્યોના જે
ગણને અનંત નિદર્શ હોય, તેને ગણનીય નિદર્શ પણ હોય છે. આ
પ્રમેય પણ પૂર્ણતા પ્રમેયનું પરિણામ છે: પૂર્ણતાની સાબિતી
વાક્યોના કોઈપણ સુસંગત ગણ માટે નિદર્શ રચે છે, અને તે નિદર્શ
ગણનીય હોય છે. ઊર્ધ્વગામી લેવેનહાઇમ--સ્કોલેમ પ્રમેય પણ છે;
તે ખાતરી આપે છે કે વાક્યોના ગણને ગણનીય નિદર્શ
હોય તો તેને અગણનીય નિદર્શ પણ હોય છે. દ્વિતીય-ક્રમ
તર્કશાસ્ત્રમાં બંને પ્રમેયો નિષ્ફળ જાય છે.
\end{explain}


\begin{thm}
\ollabel{thm:sol-no-ls} દ્વિતીય-ક્રમ તર્કશાસ્ત્રમાં લેવેનહાઇમ--સ્કોલેમ
પ્રમેય નિષ્ફળ જાય છે: એવા વાક્યો છે જેમને અનંત નિદર્શો છે,
પણ કોઈ ગણનીય નિદર્શ નથી.
\end{thm}

\begin{proof}
યાદ કરો કે
\[
\fn{Count} \ident \lexists[z][\lexists[u][\lforall[X][((X(z) \land
      \lforall[x][(X(x) \lif X(u(x)))]) \lif \lforall[x][X(x)])]]]
\]
સંરચના~$\Struct{M}$માં ત્યારે અને માત્ર ત્યારે જ સાચું છે
જ્યારે~$\Domain{M}$ ગણનીય હોય. તેથી
$\lnot\fn{Count}$ એ~$\Struct{M}$માં ત્યારે અને માત્ર ત્યારે જ
સાચું છે જ્યારે~$\Domain{M}$ અગણનીય હોય. આવી
સંરચનાઓ છે: કોઈપણ અગણનીય ગણને પ્રદેશ લો,
દાખલા તરીકે $\Pow{\Nat}$ અથવા~$\Real$. આથી
$\lnot \fn{Count}$ને અનંત નિદર્શો છે, પરંતુ કોઈ
ગણનીય નિદર્શ નથી.
\end{proof}

\begin{thm}
એવા વાક્યો છે જેમને ગણનીય નિદર્શો છે,
પણ કોઈ અગણનીય નિદર્શ નથી.
\end{thm}

\begin{proof}
$\fn{Count} \land \fn{Inf}$ એ~$\Nat$માં સાચું છે, પરંતુ
જે સંરચના~$\Struct{M}$ માટે~$\Domain{M}$
અગણનીય હોય તેમાં સાચું નથી.
\end{proof}
% END OLP-0335
\endgroup


\OLEndChapterHook
% END OLP-0330
\endgroup


\begingroup
\makeatletter\def\ol@sectioncs{section}\makeatother

% BEGIN OLP-0336
\olchapter{sol}{set}{દ્વિતીય-ક્રમ તર્કશાસ્ત્ર અને ગણસિદ્ધાંત}
\olchapterid{sol}{set}
\phantomsection\label{guunit:OLP-0336}


\begin{editorial}
આ પ્રકરણ દ્વિતીય-ક્રમ તર્કશાસ્ત્રમાં શક્તિસમૂહો અને સાતત્યકનું
સંકેતીકરણ કરે છે. પરિણામો જણાવ્યાં છે, પણ તેમની સાબિતીઓ હજી
ઉમેરવાની બાકી છે. હજી અભ્યાસપ્રશ્નો પણ નથી; વ્યાખ્યાઓ અને પરિણામોમાં
પણ ભૂલો હોઈ શકે છે. સાવચેતીથી વાંચો અને કંઈ ખોટું કે અસ્પષ્ટ લાગે
તો તેની જાણ કરો.
\end{editorial}

\begingroup
\makeatletter\def\ol@sectioncs{section}\makeatother

% BEGIN OLP-0337
\olfileid{sol}{set}{int}

\olsection{પ્રસ્તાવના}
\phantomsection\label{guunit:OLP-0337}


દ્વિતીય-ક્રમ તર્કશાસ્ત્રમાં વ્યાપકક્ષેત્રના ઉપગણો અને વિધેયો પર
પરિમાણન કરી શકાતું હોવાથી, તેમાં ઓછામાં ઓછું થોડું ગણસિદ્ધાંતનું
કામ કરી શકાય એવી અપેક્ષા સ્વાભાવિક છે. ``કામ કરી શકાય'' તેનો અહીં
અર્થ એ છે કે ગણસિદ્ધાંતના ગુણધર્મો અને વિધાનો દ્વિતીય-ક્રમ
તર્કશાસ્ત્રમાં વ્યક્ત કરી શકાય છે; તે પણ ગણો માટેની કોઈ ખાસ
અતાર્કિક શબ્દાવલી વિના (દાખલા તરીકે, ગણસિદ્ધાંતના સભ્યતા
વિધેય વિના). દાખલા તરીકે, આપણે ગણોનો સંયોગ અને છેદ તથા
ઉપગણ-સંબંધ વ્યાખ્યાયિત કરી શકીએ છીએ; ગણોનાં કદની સરખામણી પણ કરી
શકીએ છીએ અને કૅન્ટૉરના પ્રમેય જેવાં પરિણામો જણાવી શકીએ છીએ.
% END OLP-0337
\endgroup


\begingroup
\makeatletter\def\ol@sectioncs{section}\makeatother

% BEGIN OLP-0338
\olfileid{sol}{set}{cmp}

\olsection{ગણોની સરખામણી}
\phantomsection\label{guunit:OLP-0338}


\begin{prop}
સૂત્ર~$\lforall[x][(X(x) \lif Y(x))]$ ઉપગણ-સંબંધ
વ્યાખ્યાયિત કરે છે: $\Sat{M}{\lforall[x][(X(x) \lif Y(x))]}[s]$
ત્યારે અને માત્ર ત્યારે જ જ્યારે $s(X) \subseteq s(Y)$.
\end{prop}

\begin{prop}
સૂત્ર~$\lforall[x][(X(x) \liff Y(x))]$ ગણો પરનો
તાદાત્મ્ય-સંબંધ વ્યાખ્યાયિત કરે છે: $\Sat{M}{\lforall[x][(X(x)
\liff Y(x))]}[s]$ ત્યારે અને માત્ર ત્યારે જ જ્યારે $s(X) = s(Y)$.
\end{prop}

\begin{prop}
સૂત્ર~$\lexists[x][X(x)]$ અરિક્ત હોવાનો ગુણધર્મ વ્યાખ્યાયિત
કરે છે: $\Sat{M}{\lexists[x][X(x)]}[s]$ ત્યારે અને માત્ર ત્યારે જ
જ્યારે $s(X) \neq \emptyset$.
\end{prop}

ગણ~$X$ એ ગણ~$Y$ કરતાં મોટો નથી, $\cardle{X}{Y}$, ત્યારે અને માત્ર
ત્યારે જ જ્યારે કોઈ એક-એક વિધેય $f\colon X \to Y$ હોય.
આપણે વિધેય એક-એક છે અને~$X$ના ઘટકો પર તેનાં મૂલ્યો~$Y$માં છે એવું
વ્યક્ત કરી શકીએ છીએ. તેથી વ્યાપકક્ષેત્રના ઉપગણો વચ્ચે ``મોટો નથી''
સંબંધ પણ વ્યાખ્યાયિત કરી શકીએ છીએ.

\begin{prop}
નીચેનું સૂત્ર
\[
\lexists[u][(\lforall[x][(X(x) \lif Y(u(x)))] \land \lforall[x][\lforall[y][((X(x) \land X(y) \land \eq[u(x)][u(y)]) \lif
      \eq[x][y])]])]
\]
``મોટો નથી'' સંબંધ વ્યાખ્યાયિત કરે છે.
\end{prop}
\sourcecorrection{OLSOL-013}{સ્થિર અંગ્રેજી મૂળ અહીં~$u$ આખા
વ્યાપકક્ષેત્ર પર એક-એક હોવાની શરત મૂકે છે, જ્યારે સરખામણી માટે તેની
$X$ પરની મર્યાદા એક-એક હોય એટલું જ જરૂરી છે. અહીં પૂર્વપક્ષમાં
$X(x)$ અને $X(y)$ ઉમેર્યાં છે, જેથી સૂત્ર સીધું~$X$માંથી~$Y$માં
એક-એક વિધેયનું અસ્તિત્વ વ્યક્ત કરે છે.}

બે ગણો સમાન કદના અથવા ``સામ્ય'' છે, $\cardeq{X}{Y}$, ત્યારે અને
માત્ર ત્યારે જ જ્યારે કોઈ એક-એક અને વ્યાપ્ત વિધેય~$f\colon X \to Y$
હોય.

\begin{prop}
નીચેનું સૂત્ર
\begin{multline*}
  \lexists[u][(\lforall[x][(X(x) \lif Y(u(x)))] \land {}\\
    \lforall[x][\lforall[y][((X(x) \land X(y) \land \eq[u(x)][u(y)]) \lif \eq[x][y])]]
  \land {} \\\lforall[y][(Y(y) \lif \lexists[x][(X(x)
      \land \eq[y][u(x)])])])]
\end{multline*}
સામ્ય-સંબંધ વ્યાખ્યાયિત કરે છે.
\end{prop}
\sourcecorrection{OLSOL-014}{સ્થિર અંગ્રેજી મૂળનું સૂત્ર~$u$ને આખા
વ્યાપકક્ષેત્ર પર એક-એક માગે છે. પરંતુ અનંત વ્યાપકક્ષેત્રમાં જો~$X$ તેનો ઉચિત
ઉપગણ અને~$Y$ આખું વ્યાપકક્ષેત્ર હોય, તો~$X$ તથા~$Y$ સામ્ય હોવા
છતાં~$X$ પરથી~$Y$ પરનું વ્યાપ્ત એક-એક વિધેય બહારના ઘટકો સુધી
આખા વ્યાપકક્ષેત્ર પર એક-એક રીતે વિસ્તારી શકાતું નથી. અહીં
એક-એકતાની શરત~$X$ પર જ મૂકી છે.}

આ સૂત્રોને અનુક્રમે $X \subseteq Y$, $X = Y$, $X \neq
\emptyset$, $\cardle{X}{Y}$ અને $\cardeq{X}{Y}$ એમ સંક્ષેપમાં
લખીશું. (આથી થોડો ગૂંચવાડો થઈ શકે છે: ગણો~$X$ અને~$Y$ વિશે
અનૌપચારિક રીતે બોલતાં પણ આ જ સંકેતો વાપરીએ છીએ. અહીં, જોકે, આ
સંકેતો એકસ્થાની સંબંધ-ચલો~$X$ અને~$Y$ ધરાવતાં દ્વિતીય-ક્રમ
સૂત્રોના સંક્ષેપો છે.)

\begin{prop}
વાક્ય~$\lforall[X][\lforall[Y][((\cardle{X}{Y} \land
    \cardle{Y}{X}) \lif \cardeq{X}{Y})]]$ પ્રમાણભૂત છે.
\end{prop}

\begin{proof}
સંરચના~$\Struct{M}$માં આ વાક્ય ત્યારે સંતોષાય જ્યારે
કોઈપણ ઉપગણો $X \subseteq \Domain{M}$ અને $Y \subseteq
\Domain{M}$ માટે $\cardle{X}{Y}$ તથા $\cardle{Y}{X}$ હોય તો
$\cardeq{X}{Y}$ હોય. પરંતુ આ વાત \emph{કોઈપણ} ગણો~$X$ અને~$Y$
માટે સાચી છે: તે શ્રેડર--બર્નસ્ટાઇન પ્રમેય છે.
\end{proof}
% END OLP-0338
\endgroup


\begingroup
\makeatletter\def\ol@sectioncs{section}\makeatother

% BEGIN OLP-0339
\olfileid{sol}{set}{crd}

\olsection{ગણોનાં કાર્ડિનલ કદ}
\phantomsection\label{guunit:OLP-0339}


\begin{explain}
વ્યાપકક્ષેત્ર સાન્ત કે અનંત, ગણનીય કે અગણનીય
છે એવું વ્યક્ત કરી શકીએ છીએ. તેવી જ રીતે~$\Domain{M}$નો કોઈ ઉપગણ
સાન્ત કે અનંત, ગણનીય કે અગણનીય છે તે ગુણધર્મ
પણ વ્યાખ્યાયિત કરી શકીએ છીએ.
\end{explain}

\begin{prop}
સૂત્ર $\fn{Inf}(X) \ident$
\begin{multline*}
\lexists[u][(\lforall[x][\lforall[y][(\eq[u(x)][u(y)] \lif
      \eq[x][y])]] \land {} \\
\lforall[x][(X(x) \lif X(u(x)))] \land {} \\
\lexists[y][(X(y) \land \lforall[x][(X(x)
      \lif \eq/[y][u(x)])])])]
\end{multline*}
ચલ-નિયુક્તિ~$s$ને અનુલક્ષીને ત્યારે અને માત્ર ત્યારે જ સંતોષાય છે
જ્યારે~$s(X)$ અનંત હોય.
\end{prop}
\sourcecorrection{OLSOL-015}{સ્થિર અંગ્રેજી મૂળના સૂત્રમાં~$u$ને
$X$માં જ મૂલ્યો લેવાની શરત નથી, અને સૂત્રનો બહારનો કૌંસ બંધ થયો
નથી. બે ઘટકોના વ્યાપકક્ષેત્રમાં એક-ઘટક
$X$ લો અને~$u$થી બંને ઘટકો અદલો: મૂળ સૂત્ર સાચું થાય છે, છતાં~$X$
સાન્ત છે. અહીં $X(x) \lif X(u(x))$ ઉમેર્યું છે. ત્યારે~$u$ની
$X$ પરની મર્યાદા એક-એક પરંતુ~$X$ પર વ્યાપ્ત નહીં હોય ત્યારે જ
સૂત્ર સંતોષાય છે; ચૂકી ગયેલો બંધ કૌંસ પણ ઉમેર્યો છે.}

\begin{prop}
સૂત્ર $\fn{Count}(X) \ident $
\begin{multline*}
\lnot \lexists[x][X(x)] \lor {} \\
\lexists[z][\lexists[u][(X(z) \land
    \lforall[x][(X(x) \lif X(u(x)))] \land {} \\ \lforall[Y][((Y(z) \land
      \lforall[x][(Y(x) \lif Y(u(x)))]) \lif X \subseteq Y)])]]
\end{multline*}
ચલ-નિયુક્તિ~$s$ને અનુલક્ષીને ત્યારે અને માત્ર ત્યારે જ સંતોષાય છે
જ્યારે~$s(X)$ ગણનીય હોય.
\end{prop}
\sourcecorrection{OLSOL-016}{સ્થિર અંગ્રેજી મૂળ દરેક~$Y$ માટે
$X=Y$ માગે છે; પરંતુ આખું વ્યાપકક્ષેત્ર પોતે~$z$ ધરાવતું અને~$u$
હેઠળ બંધ~$Y$ છે, તેથી મૂળ સૂત્ર કોઈ ઉચિત ઉપગણ~$X$ માટે સાચું થઈ
શકતું નથી. જરૂરી નિષ્કર્ષ $X\subseteq Y$ છે: $z$ની~$u$-કક્ષા
$X$ના બધા ઘટકો સુધી પહોંચે છે. ખાલી ઉપગણ પણ પરિગણનીય ગણાય તે
માટે તેની અલગ વિસંયોજિકા ઉમેરી છે.}

કૅન્ટૉરના પ્રમેયથી આપણે જાણીએ છીએ કે અગણનીય ગણો છે;
અનંત કદનાં અનંત જુદાં સ્તરો પણ છે. ગણસિદ્ધાંત ગણોનાં કદનું આખું
અંકગણિત વિકસાવે છે અને ગણોને અનંત કાર્ડિનલ સંખ્યાઓ આપે છે. પ્રાકૃતિક
સંખ્યાઓ સાન્ત ગણોનાં કદ માપતી કાર્ડિનલ સંખ્યાઓ તરીકે કામ કરે છે.
ગણનીય ગણોનું કદ પહેલું અનંત કાર્ડિનલ કદ છે, જેને~$\aleph_0$
(``આલેફ-નૉટ'' અથવા ``આલેફ-શૂન્ય'') કહે છે. પછીનું અનંત કદ~$\aleph_1$
છે: સૌથી નાનું અપરિગણનીય કાર્ડિનલ કદ. ``$X$નું કદ~$\aleph_0$ છે''
તેને $\fn{Aleph}_0(X) \liff \fn{Inf}(X) \land \fn{Count}(X)$
થી વ્યાખ્યાયિત કરી શકીએ છીએ. $X$નું કદ~$\aleph_1$ ત્યારે અને માત્ર
ત્યારે જ જ્યારે~$X$ પોતે અપરિગણનીય હોય અને તેનો દરેક અપરિગણનીય
ઉપગણ~$X$ સાથે સામ્ય હોય. તેથી આ ગુણધર્મ સૂત્ર
$\fn{Aleph}_1(X) \ident \lnot\fn{Count}(X) \land
\lforall[Y][((Y \subseteq X \land \lnot\fn{Count}(Y))
  \lif \cardeq{X}{Y})]$ દ્વારા વ્યક્ત કરી શકાય છે. $\aleph_2$નું
કદ પણ એ જ રીતે વ્યાખ્યાયિત કરી શકાય છે; વગેરે.
\sourcecorrection{OLSOL-017}{સ્થિર અંગ્રેજી મૂળ કહે છે કે
$\aleph_1$-કદના~$X$ના બધા ઉપગણો સાન્ત અથવા~$\aleph_0$-કદના છે;
પરંતુ~$X$ પોતે પોતાનો ઉપગણ હોવાથી આ અશક્ય છે. અહીં અપરિગણનીય
ઉપગણો~$X$ સાથે સામ્ય હોવાની લઘુત્તમતા-શરત મૂકી છે, અને
\texttt{\textbackslash fn\{Aleph\_1\}}ને અગાઉના
\texttt{\textbackslash fn\{Aleph\}\_0} સાથે સુસંગત
\texttt{\textbackslash fn\{Aleph\}\_1} કર્યું છે.}

એક કદ ખાસ રસનું છે: સાતત્યકનું કાર્ડિનલ કદ. તે~$\Pow{\Nat}$નું
અથવા, સમકક્ષ રીતે, $\Real$નું કદ છે. કોઈ ગણ સાતત્યકના કદનો છે
તે દ્વિતીય-ક્રમ તર્કશાસ્ત્રમાં વ્યક્ત કરી શકાય છે, જોકે તે માટે
થોડું વધુ કામ કરવું પડે છે.
% END OLP-0339
\endgroup


\begingroup
\makeatletter\def\ol@sectioncs{section}\makeatother

% BEGIN OLP-0340
\olfileid{sol}{set}{pow}

\olsection{સાતત્યકનું કદ}
\phantomsection\label{guunit:OLP-0340}


\begin{explain}
દ્વિતીય-ક્રમ તર્કશાસ્ત્રમાં પ્રદેશના ઉપગણો પર પરિમાણન કરી
શકીએ છીએ, પરંતુ પ્રદેશના ઉપગણોના ગણો પર સીધું પરિમાણન કરી
શકતા નથી. તે માટે \emph{તૃતીય-ક્રમ} તર્કશાસ્ત્ર જોઈએ. દાખલા તરીકે,
ગણના શક્તિસમૂહમાંથી તે ગણમાં કોઈ એક-એક વિધેય નથી એવો
કૅન્ટૉરનો પ્રમેય જણાવવો હોય, તો આપણે ``દરેક ગણ~$X$ અને દરેક
ગણ~$P$ માટે, જો~$P$ એ~$X$નો શક્તિસમૂહ હોય, તો
$\cardle{P}{X}$ નહીં'' એમ લખવાનો પ્રયાસ કરી શકીએ. પરંતુ~$P$
એ~$X$નો શક્તિસમૂહ છે એમ કહેવા માટે~$P$ના ઘટકો એ~$X$ના
બધા અને માત્ર ઉપગણો છે એમ ઔપચારિક રીતે વ્યક્ત કરવું પડે; દાખલા
તરીકે $\lforall[Y][(P(Y) \liff Y \subseteq X)]$ જેવું કંઈક.
અડચણ~$P(Y)$માં છે: તે દ્વિતીય-ક્રમ તર્કશાસ્ત્રનું સૂત્ર
નથી, કારણ કે~$P$ જેવા એકસ્થાની સંબંધ-ચલોના દલીલો માત્ર પદો હોઈ
શકે છે.

છતાં, વ્યાપકક્ષેત્ર પૂરતું મોટું હોય તો ગણોના ગણો પરના પરિમાણનનું
\emph{અનુકરણ} કરી શકીએ છીએ. વિચાર એ છે કે દ્વિસ્થાની સંબંધ~$R$
પ્રદેશના ઘટકોને પ્રદેશના જ ઘટકો સાથે જોડે છે.
એવો~$R$ હોય તો કોઈ~$x$ સાથે~$R$-સંબંધિત બધા ઘટકોનો ગણ લઈએ:
$\Setabs{y \in \Domain{M}}{R(x,y)}$ એ~$x$ દ્વારા
``સંકેતિત'' ગણ છે. ઊલટું, $Z \subseteq \Pow{\Domain{M}}$
એ~$\Domain{M}$ના ઉપગણોનો કોઈ ગણ હોય અને~$\Domain{M}$માં
$Z$ના ગણોની સંખ્યા જેટલા અથવા વધુ ઘટકો હોય, તો એવો
સંબંધ~$R \subseteq \Domain{M}^2$ પણ મળે છે કે દરેક $Y \in Z$
કોઈ~$x$ દ્વારા~$R$નો ઉપયોગ કરીને સંકેતિત થાય.
\end{explain}

\begin{defn}
જો $R \subseteq \Domain{M}^2$ હોય, તો \emph{$x$ એ~$R$ વડે}
$\Setabs{y \in \Domain{M}}{R(x, y)}$ને \emph{સંકેતિત કરે છે}.
\end{defn}

જો ઘટક~$x \in \Domain{M}$ એ~$R$ વડે
$Z \subseteq \Domain{M}$ ગણને સંકેતિત કરે, તો
$Y \subseteq \Domain{M}$ ગણોના ગણને, એટલે~$Y$ના ઘટકો
દ્વારા સંકેતિત ગણોને, સંકેતિત કરે છે. તેથી ગણ~$Y$ એ~$R$ વડે
$\Pow{X}$ને સંકેતિત કરી શકે છે. એવું ત્યારે અને માત્ર ત્યારે જ બને
જ્યારે દરેક $Z \subseteq X$ માટે કોઈ $x \in Y$ એ~$R$ વડે~$Z$ને
સંકેતિત કરે અને દરેક $x \in Y$ એ~$R$ વડે કોઈ $Z \subseteq X$ને
સંકેતિત કરે.

\begin{prop}
સૂત્ર
  \[
  \fn{Codes}(x, R, Z) \ident \lforall[y][(Z(y) \liff
    R(x, y))]
  \]
એ વ્યક્ત કરે છે કે $s(x)$ એ~$s(R)$ વડે~$s(Z)$ને સંકેતિત કરે છે.
નીચેનું સૂત્ર
\begin{multline*}
  \fn{Pow}(Y, R, X) \ident \\
  \lforall[Z][(Z \subseteq X \lif \lexists[x][(Y(x) \land
    \fn{Codes}(x, R, Z))])] \land {} \\ \lforall[x][(Y(x) \lif
  \lforall[Z][(\fn{Codes}(x, R, Z) \lif Z \subseteq X)])]
\end{multline*}
એ વ્યક્ત કરે છે કે $s(Y)$ એ~$s(R)$ વડે~$s(X)$ના શક્તિસમૂહને
સંકેતિત કરે છે: $s(Y)$ના ઘટકો~$s(R)$ વડે~$s(X)$ના બધા અને માત્ર
ઉપગણોને સંકેતિત કરે છે.
\end{prop}
\sourcecorrection{OLSOL-018}{સ્થિર અંગ્રેજી મૂળના~$\fn{Pow}$
સૂત્રના બીજા સંયોજકમાં~$Y(x)$થી શરૂ થતો કૌંસ બંધ થતો નથી.
અહીં અંતે એક બંધ કૌંસ ઉમેર્યો છે, જેથી બહારનું નિહિતાર્થ અને બંને
સર્વપરિમાણકો યોગ્ય રીતે બંધ થાય.}

\begin{explain}
આ રીતથી ઉપગણો પર સીધું પરિમાણન કરવાને બદલે તેમના સંકેતો પર
પરિમાણન કરીને શક્તિસમૂહ વિશેનાં વિધાનો વ્યક્ત કરી શકીએ છીએ.
દાખલા તરીકે, કૅન્ટૉરનો પ્રમેય હવે એ રીતે કહી શકાય કે~$X$ના
શક્તિસમૂહને સંકેતિત કરતા કોઈપણ સંબંધના સંકેતોના ગણમાંથી~$X$
સુધી કોઈ એક-એક વિધેય નથી.
\end{explain}

\begin{prop}
નીચેનું વાક્ય
\begin{multline*}
  \lforall[X][\lforall[Y][\lforall[R][(\fn{Pow}(Y, R, X) \lif \\
      \lnot \lexists[u][(\lforall[x][\lforall[y][(\eq[u(x)][u(y)] \lif
            \eq[x][y])]] \land {}\\
        \lforall[x][(Y(x) \lif X(u(x)))])])]]]
\end{multline*}
પ્રમાણભૂત છે.
\end{prop}

\begin{explain}
ગણનીય ગણનો શક્તિસમૂહ અગણનીય હોય છે;
તેથી તેનું કાર્ડિનલ કદ કોઈપણ ગણનીય ગણના કદ~$\aleph_0$
કરતાં મોટું છે. $\Pow{\Nat}$ના કદને ``સાતત્યકનું કદ'' કહે છે,
કારણ કે તે વાસ્તવિક સંખ્યારેખા પરનાં બિંદુઓ, એટલે~$\Real$, જેટલું
છે. જો પ્રદેશમાં ગણનીય ગણનો શક્તિસમૂહ સંકેતિત
કરવા જેટલા ઘટકો હોય, તો $\aleph_0$-કદના ગણ~$X$નો શક્તિસમૂહ
સંકેતિત કરતો કોઈપણ ગણ~$Y$ તેના જેટલા કદનો હોય એમ કહીને કોઈ ગણ
સાતત્યકના કદનો છે તે વ્યક્ત કરી શકીએ છીએ. (વ્યાપકક્ષેત્ર એટલું
મોટું ન હોય, એટલે તેમાં~$\Real$ સાથે સામ્ય કોઈ ઉપગણ ન હોય, તો
$\Pow{X}$ને સંકેતિત કરતો સંબંધ પણ હોઈ શકતો નથી.)
\end{explain}

\begin{prop}
જો $\cardle{\Real}{\Domain{M}}$ હોય, તો નીચેનું સૂત્ર
\begin{multline*}
\fn{Cont}(Y) \ident
\lexists[X][\lexists[R][((\fn{Aleph}_0(X) \land
      \fn{Pow}(Y, R, X)) \land \\ 
      \lforall[x][\lforall[y][((Y(x) \land {}
      Y(y) \land \lforall[z][R(x,z) \liff R(y,z)]) \lif \eq[x][y])]])]]
\end{multline*}
$\cardeq{s(Y)}{\Real}$ વ્યક્ત કરે છે.
\end{prop}

\begin{proof}
$\fn{Pow}(Y, R, X)$ વ્યક્ત કરે છે કે~$s(Y)$ એ~$s(R)$ વડે
$s(X)$ના શક્તિસમૂહને સંકેતિત કરે છે; અને
$\fn{Aleph}_0(X)$ કહે છે કે તે ગણનીય અનંત છે. તેથી~$s(Y)$
ઓછામાં ઓછો સાતત્યકના કદનો છે, જોકે તે મોટો પણ હોઈ શકે છે
(જો~$s(Y)$ના ઘણા ઘટકો~$X$ના એક જ ઉપગણને સંકેતિત કરે).
છેલ્લો સંયોજક એ શક્યતા દૂર કરે છે: તે~$s(Y)$ના ઘટકો અને~$s(X)$ના
ઉપગણો વચ્ચેનો~$s(R)$ દ્વારા મળતો સંબંધ એક-એક હોય એમ માગે છે.
\sourcecorrection{OLSOL-019}{સ્થિર અંગ્રેજી મૂળની સાબિતીમાં છેલ્લે
$s(Z)$ આવે છે, પરંતુ અહીં પરિમાણિત ગણ~$X$ છે; તેથી અહીં~$s(X)$
લખ્યું છે. મૂળનું ``be'' પણ વ્યાકરણભૂલ છે.}
\end{proof}

\begin{prop}
$\cardeq{\Domain{M}}{\Real}$ ત્યારે અને માત્ર ત્યારે જ જ્યારે
\begin{multline*}
  \Sat{M}{\lexists[Y][(\fn{Cont}(Y) \land \\
          \lexists[u][(\lforall[x][\lforall[y][(\eq[u(x)][u(y)] \lif \eq[x][y])]] \land {} \\
\lforall[y][(Y(y) \liff \lexists[x][\eq[y][u(x)]])])])]}.
\end{multline*}
\end{prop}
\sourcecorrection{OLSOL-020}{સ્થિર અંગ્રેજી મૂળનું અહીંનું વાક્ય
માત્ર~$Y$ શક્તિસમૂહને સંકેતિત કરે છે અને~$Y$ એ~$u$ની છબીનો
ઉપગણ છે એમ માગે છે; ઓળખ-વિધેયથી આ બીજી શરત સહેલાઈથી મળે છે,
અટલે વાક્ય મોટા વ્યાપકક્ષેત્રોમાં પણ સાચું થઈ શકે છે. વળી~$Y$
પર સંકેતીકરણ એક-એક છે તેવી શરત નથી. અહીં~$\fn{Cont}(Y)$થી
$Y$નું કદ સાતત્યક જેટલું નિશ્ચિત કર્યું છે અને~$u$નું મૂલ્યસમૂહ
બરાબર~$Y$ રાખ્યું છે, જેથી આખા વ્યાપકક્ષેત્રનું કદ પણ એટલું જ થાય.}

\begin{explain}
સાતત્યક પરિકલ્પના એવું વિધાન છે કે સાતત્યકનું કદ પહેલું
અગણનીય કાર્ડિનલ કદ છે, એટલે~$\Pow{\Nat}$નું
કદ~$\aleph_1$ છે.
\end{explain}

\begin{prop}
સાતત્યક પરિકલ્પના સાચી છે ત્યારે અને માત્ર ત્યારે જ જ્યારે
\[\fn{CH} \ident
\lforall[X][(\fn{Aleph}_1(X) \liff \fn{Cont}(X))]\]
પ્રમાણભૂત હોય.
\end{prop}

$\lnot \fn{CH}$ પ્રમાણભૂત હોય ત્યારે અને માત્ર ત્યારે જ સાતત્યક
પરિકલ્પના ખોટી છે, એવું માનવું સાચું નથી. ગણનીય
વ્યાપકક્ષેત્રમાં~$\aleph_1$-કદના ઉપગણો પણ નથી અને સાતત્યકના
કદના ઉપગણો પણ નથી; તેથી ગણનીય વ્યાપકક્ષેત્રમાં
$\fn{CH}$ હંમેશાં સાચું છે.
પરંતુ સાતત્યક પરિકલ્પના ખોટી છે ત્યારે અને માત્ર ત્યારે જ પ્રમાણભૂત
હોય એવું બીજું વાક્ય આપી શકીએ છીએ:

\begin{prop}
સાતત્યક પરિકલ્પના ખોટી છે ત્યારે અને માત્ર ત્યારે જ
\[\fn{NCH} \ident
\lforall[X][(\fn{Cont}(X) \lif \lexists[Y][(Y \subseteq X \land \lnot
    \fn{Count}(Y) \land \lnot \cardeq{X}{Y})])]\]
પ્રમાણભૂત હોય.
\end{prop}
% END OLP-0340
\endgroup


\OLEndChapterHook
% END OLP-0336
\endgroup


\OLEndPartHook
% END OLP-0322
